\documentclass{article}

\usepackage{microtype}
\usepackage{graphicx}
\usepackage{booktabs}
\usepackage{hyperref}
\newcommand{\theHalgorithm}{\arabic{algorithm}}

\usepackage[preprint]{icml2026}

\usepackage{amsmath}
\usepackage{amssymb}
\usepackage{mathtools}
\usepackage{amsthm}
\usepackage[capitalize,noabbrev]{cleveref}

\theoremstyle{plain}
\newtheorem{theorem}{Theorem}
\newtheorem{proposition}[theorem]{Proposition}
\newtheorem{lemma}[theorem]{Lemma}
\newtheorem{corollary}[theorem]{Corollary}
\theoremstyle{remark}
\newtheorem{remark}[theorem]{Remark}

\newcommand{\E}{\mathbb{E}}
\newcommand{\KL}{\mathrm{KL}}
\newcommand{\Jh}{\hat{J}}
\newcommand{\Cov}{\mathrm{Cov}}
\newcommand{\Var}{\mathrm{Var}}

\icmltitlerunning{Goodhart in the Planner}

\begin{document}

\twocolumn[
  \icmltitle{Goodhart in the Planner: Sampling-Based Control Exploits\\Rewards Learned from Human Preferences}

  \icmlsetsymbol{equal}{*}

  \IfFileExists{authors.tex}{% Author block, kept out of the code repository; main.tex falls back to an anonymous placeholder without it.
\begin{icmlauthorlist}
  \icmlauthor{Soovadeep Bakshi}{equal,ut}
  \icmlauthor{Apurba Bose}{equal,ind}
\end{icmlauthorlist}

\icmlaffiliation{ut}{Walker Department of Mechanical Engineering, The University of Texas at Austin, Austin, TX, USA}
\icmlaffiliation{ind}{Independent researcher}

\icmlcorrespondingauthor{Soovadeep Bakshi}{soovadeep.bakshi@utexas.edu}
}{%
    \begin{icmlauthorlist}
      \icmlauthor{Anonymous Authors}{anon}
    \end{icmlauthorlist}
    \icmlaffiliation{anon}{Anonymous Institution}
    \icmlcorrespondingauthor{Anonymous Authors}{anon.email@domain.com}
  }

  \icmlkeywords{reward hacking, overoptimization, preference learning, model predictive control, MPPI, best-of-N}

  \vskip 0.3in
]

\printAffiliationsAndNotice{\icmlEqualContribution}

\begin{abstract}
Robots increasingly choose actions by searching against a reward model learned from human preferences, in
model-predictive control or best-of-$N$ selection; we measure what this does to the true objective. On
HalfCheetah and Hopper with crowd-sourced labels, and on Meta-World drawer-open and plate-slide with synthetic labels
and blind labels from the two authors, harder MPPI search cuts true return by up to 76\%: with human labels the drawer opens in 0 of 20
runs against 10 of 10 for the cloned policy the planner starts from, while held-out label accuracy stays at
0.75--0.97. On plate-slide human labels are exploited less than synthetic ones, and on window-open search helps. True return roughly follows how far the planner's weights move from
its proposals ($\sqrt{\KL}$) for both weighting rules we test, and best-of-$N$ with a verifier from synthetic labels
falls below a single noisy action. We prove that $\chi^2$-regularized MPPI at a temperature set by the scorer's error spread on the
proposals is worst-case optimal, that no positive gain can be certified once that spread exceeds the scorer's own, and
that an episode's return splits exactly into per-step terms, of which the guarantee covers the selection terms.
Calibrated from 100 proposal pairs ranked by true return (an oracle) on start states the evaluation never uses, this
rule is the safest fix we test: with each of five independent sets it succeeds in every arm run and keeps on average 86--94\% of
the cloned policy's return on plate-slide and drawer-open, though it is cautious where search helps.
Calibration from a person's comparisons cannot detect a bias the scorer learned from that person; correcting for the
pairs people cannot tell apart, or a second person, limits the damage.
\end{abstract}

\section{Introduction}
\label{sec:intro}

A robot that plans against a learned reward is an optimizer pointed at a proxy. Reward models learned from human
preferences are now used inside model-predictive control \citep{krupa2025learning,elhasnaouy2026active,nabail2026ubp2},
and robot foundation models are scaled at test time by sampling many candidate actions and executing the one a learned
verifier likes best \citep{kwok2025robomonkey}. In language-model alignment it is well established that optimizing
too hard against a learned reward lowers the quality it was meant to measure \citep{gao2023scaling,
coste2024ensembles,eisenstein2024helping}, and a line of theory explains when regularization toward a reference
policy does and does not prevent it \citep{huang2025chipo,huang2025bestofn,kwa2024catastrophic,laidlaw2025correlated}.
Robot planners differ: the optimizer is a sampling-based search whose strength is set by a temperature and a sample
count, the regularizer is implicit in the proposals, the plan is re-optimized at every step, and the labels number
hundreds to thousands, not hundreds of thousands. How much such a planner exploits its scorer has not been measured.

We study Model Predictive Path Integral control \citep[MPPI;][]{williams2017mppi,williams2017information} with proposals
around a policy cloned from the labelled data, the planning analogue of RLHF's reference model, and a Bradley--Terry
scorer trained from preference labels. We use the true simulator for planning, so that every failure is attributable to
the scorer and not to a learned dynamics model. Our contributions are:
\begin{itemize}
\item \textbf{Goodhart in the planner} (\cref{sec:exp-lambda}). On HalfCheetah and Hopper with real crowd-sourced labels
  \citep{yuan2024unirlhf}, and on Meta-World drawer-open and plate-slide with synthetic labels and blind labels from
  the two authors \citep{yu2020metaworld}, harder search lowers true return below that of the cloned policy it starts from, while
  a perfect scorer gains or, on hopper, loses less. High agreement on decided human comparisons does not make planning
  safe: on the drawer the pairs a person cannot tell apart carry what the planner needs. On plate-slide human labels
  are exploited less than synthetic ones, and on window-open search helps.
\item \textbf{An approximate collapse onto distance} (\cref{sec:exp-law}). Within each task, true return falls near one
  curve in $d=\sqrt{\KL}$ between the planner's weights and its proposals for both weighting rules we test, so $\chi^2$
  weights help mainly by moving less.
\item \textbf{How hard to search, and when not to} (\cref{sec:theory}). Against an error spread $\sigma$ on the
  proposals, $\chi^2$-regularized MPPI at $\lambda\sqrt{\chi^2}=\sigma$ is worst-case optimal, no positive gain can be
  certified once $\sigma$ exceeds the scorer's own spread, and the guarantee of hard exponential search falls as
  $-\sigma\sqrt N$; an exact decomposition of an episode's return shows which terms this controls. The mathematics is elementary;
  what is new is bounding the error on the planner's own proposals, which held-out accuracy does not predict
  (\cref{sec:exp-mech}), and following it through the closed loop.
\item \textbf{Fixes, measured honestly} (\cref{sec:exp-fix}). Calibrated from 100 proposal pairs ranked by true return,
  the worst-case-optimal rule is the safest fix we test, though cautious where search helps; calibrated from the
  labeller's own comparisons, it cannot detect a bias the scorer shares with that person; ensembles, $\chi^2$ weights, automatic temperatures
  and a guard each win somewhere and lose elsewhere, and a UBP2-inspired CEM planner \citep{nabail2026ubp2} is the most
  exploited. Labelling the planner's own trajectories reduces the exploitation.
\item \textbf{Best-of-$N$ is not safe either} (\cref{sec:exp-bon}). With a verifier from 500 synthetic pairs it falls below
  one noisy action from $N=16$ ($N=4$ on plate-slide); with 2,000 pairs selection helps.
\end{itemize}

\section{Setting}
\label{sec:setting}

\paragraph{Planner.} At each control step MPPI rolls out $N$ action sequences of length $H$ from the current state and
executes the first action of their weighted average \citep{williams2017mppi}. We draw the sequences around the cloned
policy's own closed-loop plan $U$: $A_i=\mathrm{clip}(U+\epsilon_i)$, $\epsilon_i\sim\mathcal N(0,\sigma^2I)$, which
plays the role of the reference policy in RLHF. Each plan gets a learned return $\Jh(A_i)=\sum_{h<H}\hat r(o_h,a_h)$
and, for evaluation only, a true return $J(A_i)$. Standard (exponential) MPPI uses
$w_i\propto\exp(\Jh_i/\lambda)$, the solution of $\max_w \sum_i w_i\Jh_i-\lambda\KL(w\Vert\mathrm{unif})$, the same
KL-regularized objective as RLHF with $\beta=\lambda$. The $\chi^2$ rule of Tsallis-divergence MPPI
\citep{wang2021tsallis} uses $w_i\propto\max(0,\Jh_i-c)$ with $c$ set by $\lambda$, the solution with $\lambda\chi^2/2$ in
place of $\lambda\KL$. Lower $\lambda$ means harder search; $\lambda\to0$ is best-of-$N$ over plans. We log the
per-step $\KL(w\Vert\mathrm{unif})=\sum_iw_i\log(Nw_i)\le\log N$. Uniform weights ($\lambda=10^6$) give the planner
without search, which we call \emph{uniform averaging}: it executes the average of the clipped noisy proposals, which
is not the cloned policy's own action, so we also run the cloned policy directly. Unless stated, $N=256$, $H=20$,
$\sigma=0.5$, and planning uses the true MuJoCo dynamics \citep{todorov2012mujoco}.

\paragraph{Scorers.} A two-layer MLP $\hat r(o,a)$ is trained on segment pairs with the Bradley--Terry likelihood
$P(\sigma^1\succ\sigma^2)=\mathrm{sigmoid}(\sum_t\hat r(\sigma^1_t)-\sum_t\hat r(\sigma^2_t))$
\citep{bradley1952rank,christiano2017deep}, with 10\% of pairs held out, and standardized per step; a tie (``can't
tell'') is a target of $\frac12$, that is, indifference. Five seeds per
label set give single scorers and a five-member ensemble \citep{lakshminarayanan2017ensembles}.

\paragraph{Tasks and labels.} \emph{Locomotion}: HalfCheetah-v5 and Hopper-v5 \citep{towers2024gymnasium}, D4RL
medium-replay data \citep{fu2020d4rl} and the Uni-RLHF crowd-sourced labels \citep{yuan2024unirlhf}: 2,000 pairs of
200-step segments that agree with the true ranking on 73\% (cheetah) and 64\% (hopper) of pairs, with 10--15\% ties.
\emph{Manipulation}: Meta-World v3 drawer-open, plate-slide and window-open \citep{yu2020metaworld} with mocap position control.
The cloned policy is trained on 10 scripted demonstrations; the labelled data add 100 rollouts of the cloned policy
with action noise from 0 to 1, and 2,000 pairs of 50-step segments are labelled by a noisy Bradley--Terry teacher that
agrees with the true ranking on 86--88\% of pairs. Smaller budgets use the first $K$ pairs. The authors also labelled
pairs blind (two clips side by side in random order, with ``can't tell''): drawer pairs 0--299 and 200--499 (97.6\%
and 97.7\% agreement with the true ranking on decided pairs, 18\% and 13\% ties; the two agreed on all 73 shared
decided pairs) and plate-slide pairs 0--299 (97.2\%, 17\%). Episodes last 200 steps on
the walkers and 150 on the arm. True return is the environment's reward (Meta-World's dense reward on the arm), and on
the arm we also count successes. Each cell averages 10 runs (5 scorer seeds $\times$ 2 start states; 10 start states for
ensembles), and 20 for standard MPPI in the five main settings. Appendix~\ref{app:setup} lists all settings; code and results will be released.

\section{Theory: what the planner can and cannot certify}
\label{sec:theory}

Fix a state. Let $q$ give each of the $N$ proposals weight $1/N$ (the reference policy of RLHF), $J$ and $\Jh$ be the
true and learned plan returns, and $e=\Jh-J$ the scorer's error; labels fix $\Jh$ only up to scale and offset, so $J$
is measured in the scorer's units. A plan distribution $p\ll q$ is any reweighting of the proposals, such as MPPI's
weights; its density $w=dp/dq$ gives plan $i$ the factor $w_i=Np_i$ by which it is up-weighted ($w_i=2$ doubles its
share, $w_i=0$ drops it). Let $\Delta(p)=\E_p[J]-\E_q[J]$ be the true gain of selecting with $p$, $\hat\Delta(p)$ the
same under $\Jh$, and $s(p)=\sqrt{\chi^2(p\Vert q)}=\big(\tfrac1N\sum_i(w_i-1)^2\big)^{1/2}$ the distance moved: $s=0$
is no search, and spreading the weight evenly over the best $k$ plans gives $s=\sqrt{N/k-1}$ (1 for the best half,
$\approx16$ for the single best of 256). What matters is $\sigma_e=\mathrm{sd}_q(e)$, the spread of the scorer's mistake
across the proposals: an error shared by all of them changes no weights and costs nothing; what hurts is overrating
some proposals relative to others. Proofs are in Appendix~\ref{app:proofs}.

\begin{lemma}[Tight change of measure]\label{lem:com}
For every $p\ll q$, $|\Delta(p)-\hat\Delta(p)|\le\sigma_e\,s(p)$. For each $p$, an error of spread $\sigma_e$ that
overrates every plan in proportion to its extra weight $w-1$ attains $\Delta(p)=\hat\Delta(p)-\sigma_e s(p)$.
\end{lemma}

In words, the gain the scorer reports overstates the true gain by at most how far the planner moves times how much
the scorer's mistakes vary across the proposals, and by exactly that much when the scorer overrates precisely the plans
the planner up-weights. This sharpens the change-of-measure bound used for language models
\citep{huang2025chipo,huang2025bestofn}: it vanishes when the planner does not move, and no smaller bound holds. Against every error of spread at most $\sigma$, the
planner can therefore guarantee exactly $W_\sigma(p)=\hat\Delta(p)-\sigma s(p)$, and the best it can do is maximize this.

\begin{theorem}[Worst-case-optimal search]\label{thm:chi2}
(i) If $\sigma\ge\sigma_q(\Jh)$, no plan distribution guarantees a gain: $\max_pW_\sigma(p)=0$, attained by $p=q$.
\\[2pt](ii) If $\sigma<\sigma_q(\Jh)$, $W_\sigma$ is maximized by $\chi^2$ MPPI, $w=\max(0,(\Jh-c)/\lambda)$, at a
temperature $\lambda^\star$ with $\lambda^\star s(p_{\lambda^\star})=\sigma$ (for $\sigma=0$, the hard-search limit
$\lambda\to0$), and the maximum is positive. Along the
$\chi^2$ path, $\lambda\,s(p_\lambda)$ falls from $\sigma_q(\Jh)$ to $0$ as search hardens, and $W_\sigma$ rises until
$\lambda^\star$ and falls after it.
\\[2pt](iii) While no proposal's weight is clipped to zero, $\hat\Delta(p_\lambda)=\sigma_q(\Jh)\,s$ and the true gain
is exactly $\Delta(p_\lambda)=\rho_q\,\sigma_q(J)\,s$, where $\rho_q=\mathrm{Corr}_q(J,\Jh)$.
\end{theorem}

The theorem turns the error size into a stopping rule the planner can check from its own weights: harden the search
until $\lambda\sqrt{\chi^2}$ falls to the error spread, and do not search at all when the error spread exceeds the
scorer's own spread over the proposals. The $\chi^2$ geometry comes from the error model, not a choice of divergence:
knowing only the error's spread makes $\sigma_e\sqrt{\chi^2}$ the exact worst case (\cref{lem:com}), while sub-Gaussian
errors would give a bound in $\sqrt{\KL}$ matched by exponential MPPI, an assumption heavy tails break
\citep{kwa2024catastrophic}.

\paragraph{Scale.} Since $J$ is known only up to scale, we put it in the scorer's units by the rescaling that best
matches $\Jh$ on the proposals (least squares of $\Jh$ on $J$), which leaves an error uncorrelated with $J$ and gives
$\sigma_e=\sigma_q(\Jh)\sqrt{1-\rho_q^2}$. The threshold of (i) then becomes $\rho_q\le0$: a
positive gain can be certified exactly when the scorer's Pearson correlation with the true return over the proposals is
positive. For jointly Gaussian returns this is the same as ordering a random pair correctly more often than chance; in
general it is not, since a scorer can order most pairs correctly and still correlate negatively by badly underrating
one high-return plan. As $\rho_q\to0^+$ the rescaling shrinks $J$, so the certified gain is positive but measured in ever smaller units: the
theorem then guarantees the sign of the gain, not its size. In scale-free form, $\chi^2$ MPPI with $\lambda s=\sigma$
does not lose for any true return whose correlation with $\Jh$ over the proposals is at least
$\sqrt{1-\sigma^2/\sigma_q(\Jh)^2}$.

MPPI executes the weighted average of the plans, $\bar A_p=\E_p[A]$, not a plan drawn from $p$. The guarantee carries
over up to the spread of the plans the weights keep.

\begin{proposition}[The executed plan]\label{prop:exec}
If $J$ is $L$-smooth in the plan on the action box, then $J(\bar A_p)\ge\E_p[J]-\tfrac L2\,\E_p\Vert A-\bar A_p\Vert^2$, so
against every error of spread at most $\sigma$, $J(\bar A_p)-\E_q[J]\ge W_\sigma(p)-\tfrac L2\,\E_p\Vert A-\bar A_p\Vert^2$.
The last term vanishes as the weights concentrate on one plan, and it can be dropped when $J$ is concave in the plan.
\end{proposition}

$L$ is unknown for contact-rich tasks, so we measure the averaging term $J(\bar A_p)-\E_p[J]$ directly
(\cref{app:probe}): at least $-0.05$ standard deviations except on window-open ($-0.15$).

\begin{corollary}[Hard exponential search]\label{prop:exp}
On $N$ samples with a unique top-scored plan, exponential MPPI with $\lambda\to0$ selects it, $s\to\sqrt{N-1}$, and its
guarantee is $\max\Jh-\E_q\Jh-\sigma\sqrt{N-1}$: the worst error overrates the selected plan by $\sigma\sqrt{N-1}$. If
the top score's margin grows more slowly than $\sqrt N$, as $\sigma_q(\Jh)\sqrt{2\log N}$ does for Gaussian scores,
this tends to $-\infty$ as $N$ grows, for every $\sigma>0$; the guarantee of \cref{thm:chi2} is never negative.
\end{corollary}

How large is $\sigma_e$? For a linear reward its bound shrinks with the label count at a rate set by how far the planner's
proposals vary beyond the labelled pairs (\cref{cor:labels}, \cref{app:proofs}). The next result locates where along the exponential path true return turns down.

\begin{proposition}[Where search starts to hurt]\label{prop:onset}
Let $p_\eta\propto q\,e^{\eta\Jh}$ with $\eta=1/\lambda$, and $d(\eta)=\sqrt{\KL(p_\eta\Vert q)}$.
(i) For every $\eta$, $\tfrac{d}{d\eta}\E_{p_\eta}[J]=\Cov_{p_\eta}(J,\Jh)$ and
$\tfrac{d}{d\eta}\KL(p_\eta\Vert q)=\eta\Var_{p_\eta}(\Jh)$. The true return of a plan drawn from the planner's
weights rises exactly while learned and true returns are positively correlated under those weights, and it has a
local peak only where that covariance changes sign from positive to negative.
\\[2pt](ii) As $\eta\to0$, $\E_{p_\eta}[J]-\E_q[J]=\sqrt2\,\rho_q\,\sigma_q(J)\,d+O(d^2)$, where
$\rho_q=\mathrm{Corr}_q(J,\Jh)$ and $\sigma_q(J)$ is the spread of true return over the proposals.
\\[2pt](iii) The executed plan $\bar A=\E_{p}[A]$ satisfies
$J(\bar A)-J(U)=\underbrace{\E_p[J]-\E_q[J]}_{\text{selection gain}}-\underbrace{(J(U)-\E_q[J])}_{\text{price of
noise}}+\underbrace{(J(\bar A)-\E_p[J])}_{\text{averaging}}$; on $N$ samples the last term vanishes as
$\lambda\to0$, so hard search beats the cloned plan only if the selection gain exceeds a price that does not
depend on the scorer.
\end{proposition}

Part (ii) makes $\sqrt{\KL}$ the natural axis, as \citet{gao2023scaling} found for language models: to first order
true return changes linearly in it, with a slope set by how well the scorer ranks the proposals, not how well it fits
its labels (near $q$, $\sqrt{\KL}\approx s/\sqrt2$, matching \cref{thm:chi2}(iii)). Part (i) says
that exploitation begins where the scorer stops ranking the planner's own favourites correctly, possibly long before
it ranks typical proposals badly. Part (iii) separates exploitation from a cost of sampling that a perfect scorer also
pays; drawing one plan from the weights, as inference-time pessimism does \citep{huang2025bestofn}, gives up the
averaging term at every temperature, which is why it does worse than averaging on cheetah (167--265) and the drawer.

Over an episode the planner re-plans at every state it visits, and the per-step terms add up exactly.

\begin{proposition}[Closed loop]\label{prop:loop}
Let the dynamics be deterministic, $V_t$ the cloned policy's true return from step $t$ to the end of the episode, and
$Q_t(s,a)$ the return of taking $a$ at step $t$ and following the cloned policy after it. If at each visited state
$s_t$ the planner executes the first action $a_t$ of $\bar A_t$, its episode return $R$ satisfies
$R-V_0(s_0)=\sum_{t<T}\big[\Delta_t-c_t+\alpha_t+\eta_t\big]$, where $\Delta_t$, $c_t$ and $\alpha_t$ are the
selection gain, price of noise and averaging term of \cref{prop:onset}(iii) at $s_t$, and
$\eta_t=Q_t(s_t,a_t)-V_t(s_t)-\big(J(\bar A_t)-J(U_t)\big)$ is the horizon gap, by which the $H$-step comparison with
the cloned plan misstates the long-run advantage of the first action. If at every visited state the weights are
worst-case optimal for an error spread of at least that state's $\sigma_e$, then every $\Delta_t\ge0$ and
$R\ge V_0(s_0)-\sum_t(c_t-\alpha_t-\eta_t)$.
\end{proposition}

The identity does not make the other terms independent of the scorer, which sets the weights and with them the
averaging term and the states visited: with $J(a)=1-a^2$, cloned action 0 and proposals $\pm1$, hard search on a
scorer that favours $+1$ has zero selection gain and loses only the averaging term. \Cref{thm:chi2} controls the
selection terms, given the error spread on the visited states, which can exceed that on the cloned policy's; we measure
the rest (\cref{app:loop}).

\section{Experiments}
\label{sec:exp}

\subsection{Harder search lowers true return}
\label{sec:exp-lambda}

\begin{figure*}[t]
\centering
\includegraphics[width=\textwidth]{figures/fig_lambda.pdf}
\vspace{-6mm}
\caption{True return against MPPI temperature (harder search to the right), mean of 10--20 runs with $\pm$1 standard
error over runs. Dotted: uniform averaging (no search); dash-dot: the cloned policy run directly. A perfect scorer (the true reward) keeps or improves its return except on hopper (\cref{sec:exp-law}); a
learned scorer loses it, except on window-open, where it gains less. Magenta: 300 blind human labels, exploited faster
than 500 synthetic ones; grey: 2,000 labels, safe on the drawer and window-open but not on plate-slide.}
\label{fig:lambda}
\end{figure*}

\Cref{fig:lambda} shows the main effect. With the crowd scorer, cheetah's true return falls from 480 at $\lambda=3$ to
262 at $\lambda=0.01$ (uniform averaging scores 233, the cloned policy run directly 412), and hopper's from 468 to 272--356, while the perfect scorer
improves on cheetah (867 to 1,010). On hopper hard search costs return even with the perfect scorer (559 at
$\lambda=3$, 342--373 at $\lambda\le1$), so part of the hopper loss is the planner's own (\cref{sec:exp-law}). On the drawer with 500 labels, true return falls from 707 to 504 and the
drawer stays shut in 9 of 20 runs, while the scorer's own estimate of the executed episode rises twelve-fold (56 to
674; uniform averaging and the cloned policy open it every time). Blind human
labels on 300 pairs are the starkest case: the gentlest search matches uniform averaging (851, 20 of 20), and every
temperature at or below 0.1 opens the drawer in 0 of 20 runs (201--216); both authors' labels on 500 pairs fare no
better (246 at $\lambda=0.01$, 0/10). On plate-slide, where the cloned policy already succeeds every time, every scorer
from synthetic labels loses ground under harder search, with 500 labels (950 to 436) and with 2,000 (977 to 319).
Human labels lose less: one author's 300 keep 1,066 (10/10) at hard search, against 499 (8/10) with synthetic labels
on the same pairs and 886 (8/10) with them in true-return order, and both authors' 500 fall to 681 (6/10).
On the drawer, 2,000 labels are enough: harder search then helps (899 to 948, never failing).
Search is not always exploited: on window-open it helps at every budget (295 without search, 375--392 at
$\lambda=0.01$; \cref{sec:exp-mech}), and button-press is a null
case (\cref{app:null}). Hard search still loses with horizon 10 or 40 and proposal noise 0.25 or 1, except on
cheetah with noise 0.25 (\cref{tab:sens}).

The human-label result is not a labelling accident. Those labels agree with the true ranking on 97.6\% of decided
pairs, against 86\% for the synthetic teacher on the same pairs, and give the best held-out accuracy (0.97) and
per-step correlation with the true reward (0.91) of any scorer we trained. High agreement on decided pairs does not make a safer planning target. Controls on the same
300 pairs point to the 55 pairs the person could not tell apart. At $\lambda=0.01$, ordering every pair by true return
gives 640 (6/10), and so does dropping 55 random pairs (672, 6/10); but with only the decided pairs in true order the
drawer never opens, whether the ties are kept as indifference (268) or dropped (233; 190 with the person's own
answers). What the planner needs is the order of those close pairs, half of the 63 in which neither clip opens the
drawer. In recorded episodes (\cref{fig:film}) the arm parks in an unvisited corner the
scorer rates very highly (synthetic labels) or pulls about 10\,cm short of the handle (human labels). A linear fit of each scorer's per-step output on four state features shows why: the human-label
scorer barely penalizes the hand--handle distance (standardized weight $-0.16$, against $-0.40$ to $-0.49$ with
synthetic labels) and weights pulling (0.29) and the drawer's opening (0.63) most.

\begin{figure}[t]
\centering
\includegraphics[width=\columnwidth]{figures/fig_filmstrip.pdf}
\caption{One drawer start under three planners (true returns). Uniform averaging opens the drawer. Standard MPPI
($\lambda=0.01$) on the 500-label scorer lifts the hand into a part of the workspace the data never visited; on the
human-label scorer ($\lambda=0.1$) it pulls in front of the handle without touching it.}
\label{fig:film}
\end{figure}

\subsection{True return follows the distance moved}
\label{sec:exp-law}

\begin{figure*}[t]
\centering
\includegraphics[width=\textwidth]{figures/fig_kl.pdf}
\vspace{-6mm}
\caption{(a) Return gained over uniform averaging, divided by the perfect scorer's best, against the
per-step distance $\sqrt{\KL}$ between MPPI's weights and its proposals. Circles: exponential weights; squares:
$\chi^2$ weights; each point is one temperature (10--20 runs); curves are least-squares fits of
$a\,d-b\,d^2$; faint dashed: perfect scorer, whose highest point is $1-{}$(averaging's return $\div$ its best), e.g.\
0.77 on cheetah and 0.08 on plate-slide. Arm tasks use 500 labels. (b) The same on the drawer for five label budgets and the
human labels. (c) Best-of-$N$ with a learned verifier on the drawer (solid) and plate-slide (dashed), relative to the
clean cloned action (dash-dot line); $N=1$ (open markers) is one noisy action. Every point is 50 runs on the same ten starts (10 for the ensemble).}
\label{fig:kl}
\end{figure*}

\Cref{fig:kl}a plots every single-scorer MPPI cell against $d=\sqrt{\KL}$. A two-parameter curve
$\Delta J=a\,d-b\,d^2$ through uniform averaging ($d=0$) fits each task well ($R^2=0.79$--$0.99$ over 7--11 cells), and
the sign of the initial slope $a$ is clear on every task (90\% hierarchical bootstrap intervals, resampling scorer
seeds and then runs: cheetah 306 to 664, hopper 163 to 330, drawer $-385$ to $-99$, plate-slide $-479$ to $-218$,
window-open 6 to 112). A curve fitted on the exponential cells alone predicts the $\chi^2$ cells with a mean absolute
error of 38\% of their spread on cheetah, 57\% on hopper, 23\% on the drawer (16\% with human labels), 22\% on
plate-slide and 29\% on window-open: at equal distance moved the two rules give similar true return, an approximate
collapse within each task for the two rules we test, not a law. The
advantage of $\chi^2$ weights at a fixed temperature (\cref{sec:exp-fix}) is therefore mostly that they move less. The
walkers first gain and then lose (peaks near $d=1.2$ and 1.0); the drawer and plate-slide, where the cloned policy is
near the optimum, lose from the start, and window-open gains throughout.

The dashed curves (the perfect scorer at the same distances) separate exploitation from the planner's own cost: on
cheetah, the drawer and plate-slide the gap opens from the first step and reaches 590--750 points at hard search, and
on window-open 175. On hopper the crowd scorer is ahead of the true reward up to $d=1.7$, possibly because its 200-step
segments reward progress the 20-step return misses, and behind it beyond $d=1.9$ ($-17$ to $-101$).

On the drawer (\cref{fig:kl}b) the fitted initial slope $a$ grows steadily with the label count: $-507$, $-382$,
$-250$, $-2$ and $+99$ for 200, 300, 500, 1,000 and 2,000 labels, roughly linear in $\log n$: $+630$ per decade
(90\% hierarchical interval $+514$ to $+766$), crossing zero near 1,200 labels (1,010--1,420). The human labels
sit below the synthetic labels on the same 300 pairs on average ($a=-460$, interval $-642$ to $-236$, against $-382$),
though the intervals overlap. With five budgets, we read the log-linear form as a summary, not a law to
extrapolate. On plate-slide, 2,000 labels do not change the sign ($a=-342$ at 500, $-360$ at 2,000).

\subsection{Two ways to be exploited, and why held-out checks do not warn}
\label{sec:exp-mech}

\begin{figure*}[t]
\centering
\includegraphics[width=\textwidth]{figures/fig_mechanism.pdf}
\vspace{-6mm}
\caption{(a) Drawer, true return against label budget at gentle and hard search; the cautious ensemble rescues 500 and
1,000 labels but not 200 or the human labels. Green: 500 labels, half on the planner's own trajectories (\cref{sec:exp-fix}). (b) Change in the scorer's own per-step estimate of the executed episode
between $\lambda=3$ and $\lambda=0.01$, against the ratio of true returns: the estimate rises (right), consistent with
extremal exploitation, or falls (left), consistent with noise chasing; diamonds: human labels. (c) Each fix at
$\lambda=0.01$, relative to standard MPPI at $\lambda=3$.}
\label{fig:mech}
\end{figure*}

\Cref{fig:mech}b separates two patterns by the sign of the change in the scorer's own estimate of the episode the
robot actually executes. On the drawer with synthetic labels the estimate rises sharply as true return falls,
consistent with the planner finding states the scorer overrates, the extremal case of \citet{manheim2018categorizing}.
On plate-slide and with human labels the estimate \emph{falls} (from 27--43 to between $-43$ and $-90$), consistent
with chasing plans that look good at planning time and do not hold up once the planner re-plans, the optimizer's curse
\citep{smith2006optimizer} in closed loop. In the scorer's observations hard search leaves the labelled data on all three (51--60\% of visited states, none
without search), but leaving it is not sufficient: on window-open 60\% of states are off it even without search, and
search still helps (\cref{app:loop}). The two patterns respond to
different fixes (\cref{sec:exp-fix}).

Why do the usual checks not warn? \Cref{tab:probe} (\cref{app:probe}) compares them with what \cref{prop:onset} says
matters, on the planner's 256 proposals along the cloned policy's episodes, scored by the learned scorer and the true
reward. The initial slope predicted by \cref{prop:onset}(ii) matches the measured one to within 6\% in eleven of the
thirteen settings (24\% and 12\%, about one standard error, with 200 drawer labels and with human labels). The correlation
$\rho_q$ that sets it is unrelated to the held-out checks: the 500-label drawer scorer has held-out accuracy 0.88 and
per-step correlation 0.88 on the labelled data, but ranks the planner's proposals with $\rho_q=0.10$; the hopper scorer
has a per-step correlation of only 0.08 on the data, yet ranks proposals with $\rho_q=0.76$. Part (iii) explains why hard
search loses to the cloned plan even when selection still helps: perturbing the cloned plan costs 0.7--2.2 standard
deviations of true plan return, and the scorers recover only a fraction of that by selection. On window-open it costs
only 0.33 and selection recovers 0.7--1.0, which is why search helps there. A scorer can be checked
only where it will be optimized: on the planner's own proposals and, after the first step, on states the cloned policy
never visited. Whole episodes bear this out (\cref{prop:loop}; \cref{app:loop}). On the states hard search visits, the
scorer ranks the proposals worse than on the cloned policy's states ($\rho$ from 0.10 to $-0.11$ on the drawer,
$-0.10$ to $-0.31$ with human labels, 0.45 to $-0.07$ on plate-slide, 0.76 to 0.51 on hopper; near zero on both on
cheetah) and stays positive on window-open (0.35 to 0.22). Summed over a drawer episode, the selection gains are $-151$ with 500
labels and $-373$ with human labels, against $+632$ with the perfect scorer. On plate-slide they stay positive ($+113$)
and the loss is in the horizon gap ($-497$; $-237$ for the perfect scorer on the drawer): the executed plans' 20-step
advantage over the cloned plan overstates their long-run effect.

The same proposals test \cref{thm:chi2}. Given the true error spread, the worst-case-optimal rule does not lose at any of
the 2,042 states and gains 0.05--0.53 standard deviations on average; standard MPPI at hard search gains more on average
in ten of thirteen settings but loses 1.6--8.0 at its worst state, which a closed-loop planner has many chances to meet.
The unknown is the error spread, which the ensemble's disagreement underestimates by a factor $\kappa^\star$
of 1.0--2.8, most with human labels: scorers trained on the same labels share their mistakes. A practical version
calibrates it from pairs of the planner's own proposals ranked by true return (an oracle; in deployment, a person's preferences):
$\sigma=\sigma_q(\Jh)\sqrt{1-\hat\rho_+^2}$, with $\hat\rho_+=\max(0,\hat\rho)$ and $\hat\rho=\sin(\pi(c-\frac12))$
from the share $c$ of pairs the scorer orders correctly \citep{kruskal1958ordinal}, so no search when $\hat\rho\le0$. With 100 pairs it loses at most
0.22 standard deviations at its worst state in twelve of thirteen settings and keeps 47--93\% of the true-spread
rule's gain where that gain exceeds 0.2; a nominal 95\% lower bound on $c$ cuts the loss to 0.05 but keeps only
12--77\% of the gain. By \cref{cor:calib} (\cref{app:calib}) the rule cannot lose at a state whose correlation is at least
$\hat\rho$; mistakes the labeller shares with the scorer, sampling noise and pooling (19--59\% of states fall below the
pooled value) can each put $\hat\rho$ above a state's own, so the rule is an empirical safeguard, not a per-state
certificate.

\subsection{What helps, and where}
\label{sec:exp-fix}

\begin{table*}[t]
\centering
\caption{True return at hard search, mean $\pm$ standard error over 10 runs per cell ($^\dagger$\,20, $^\ast$\,5); successes over runs in parentheses on the arm. $^\ddagger$\,Mean $\pm$ standard deviation across five independent 100-pair calibration sets ranked by true return (an oracle), drawn on start states the evaluation never uses, each calibrating its own controller (\cref{tab:calibcl}). Bold: best learned-scorer method per column (the calibrated temperature on the arm, the automatic one on the walkers).}
\label{tab:fixes}
\small
\begin{tabular}{lccccc}
\toprule
& Cheetah & Hopper & Drawer 500 & Drawer human & Plate 500 \\
\midrule
Cloned policy, run directly & 412\,{\scriptsize$\pm$114} & 502\,{\scriptsize$\pm$32} & 898\,{\scriptsize$\pm$1}\,(10/10) & 898\,{\scriptsize$\pm$1}\,(10/10) & 1111\,{\scriptsize$\pm$5}\,(10/10) \\
Uniform averaging (no search) & 233\,{\scriptsize$\pm$78} & 418\,{\scriptsize$\pm$38} & 852\,{\scriptsize$\pm$1}\,(10/10) & 852\,{\scriptsize$\pm$1}\,(10/10) & 1051\,{\scriptsize$\pm$6}\,(10/10) \\
Perfect scorer, $\lambda=0.01$ & 1010\,{\scriptsize$\pm$17} & 373\,{\scriptsize$\pm$34} & 1090\,{\scriptsize$\pm$7}\,(10/10) & 1090\,{\scriptsize$\pm$7}\,(10/10) & 1139\,{\scriptsize$\pm$7}\,(10/10) \\
Learned scorer, gentle ($\lambda=3$) & 480\,{\scriptsize$\pm$38}$^\dagger$ & 468\,{\scriptsize$\pm$31}$^\dagger$ & 707\,{\scriptsize$\pm$41}$^\dagger$\,(18/20) & 851\,{\scriptsize$\pm$4}$^\dagger$\,(20/20) & 950\,{\scriptsize$\pm$21}$^\dagger$\,(20/20) \\
\midrule
\multicolumn{6}{l}{\emph{Learned scorer, hard search ($\lambda = 0.01$)}} \\
Standard (exp.\ weights) & 262\,{\scriptsize$\pm$25}$^\dagger$ & 356\,{\scriptsize$\pm$39}$^\dagger$ & 504\,{\scriptsize$\pm$74}$^\dagger$\,(11/20) & 216\,{\scriptsize$\pm$10}$^\dagger$\,(0/20) & 436\,{\scriptsize$\pm$65}$^\dagger$\,(15/20) \\
$\chi^2$ weights & 364\,{\scriptsize$\pm$23} & 303\,{\scriptsize$\pm$49} & 480\,{\scriptsize$\pm$96}\,(5/10) & 211\,{\scriptsize$\pm$9}\,(0/10) & 448\,{\scriptsize$\pm$80}\,(7/10) \\
Cautious ensemble (mean $-$ sd) & 295\,{\scriptsize$\pm$30} & 374\,{\scriptsize$\pm$30} & 801\,{\scriptsize$\pm$51}\,(10/10) & 261\,{\scriptsize$\pm$42}\,(1/10) & 406\,{\scriptsize$\pm$68}\,(8/10) \\
$\chi^2$ on ensemble mean & 485\,{\scriptsize$\pm$25} & 448\,{\scriptsize$\pm$56} & 538\,{\scriptsize$\pm$85}\,(6/10) & 219\,{\scriptsize$\pm$13}\,(0/10) & 360\,{\scriptsize$\pm$75}\,(7/10) \\
Automatic guard & 409\,{\scriptsize$\pm$31} & 361\,{\scriptsize$\pm$60} & 729\,{\scriptsize$\pm$77}\,(9/10) & 222\,{\scriptsize$\pm$17}\,(0/10) & 335\,{\scriptsize$\pm$72}\,(6/10) \\
\midrule
\multicolumn{6}{l}{\emph{Learned scorer, adaptive temperature or a different optimizer}} \\
Automatic $\lambda$ ($\chi^2$, $c=0.5$) & \textbf{521}\,{\scriptsize$\pm$20} & \textbf{569}\,{\scriptsize$\pm$1} & 650\,{\scriptsize$\pm$60}\,(7/10) & 262\,{\scriptsize$\pm$23}\,(0/10) & 662\,{\scriptsize$\pm$61}\,(10/10) \\
Worst-case-optimal $\lambda$ (Thm.~\ref{thm:chi2}, $\kappa=1$) & 460\,{\scriptsize$\pm$52} & 569\,{\scriptsize$\pm$2} & 575\,{\scriptsize$\pm$76}\,(8/10) & 231\,{\scriptsize$\pm$13}\,(0/10) & 699\,{\scriptsize$\pm$58}\,(10/10) \\
Worst-case-optimal $\lambda$ (Thm.~\ref{thm:chi2}, $\kappa=2$) & 253\,{\scriptsize$\pm$84} & 426\,{\scriptsize$\pm$32} & 821\,{\scriptsize$\pm$8}\,(10/10) & 382\,{\scriptsize$\pm$75}\,(2/10) & 883\,{\scriptsize$\pm$47}\,(10/10) \\
Worst-case-optimal $\lambda$, 100 oracle pairs & 357\,{\scriptsize$\pm$86}$^\ddagger$ & 539\,{\scriptsize$\pm$12}$^\ddagger$ & \textbf{840}\,{\scriptsize$\pm$16}$^\ddagger$\,(50/50) & \textbf{852}\,{\scriptsize$\pm$0}$^\ddagger$\,(50/50) & \textbf{959}\,{\scriptsize$\pm$86}$^\ddagger$\,(50/50) \\
CEM, cautious & 472\,{\scriptsize$\pm$9} & 337\,{\scriptsize$\pm$66} & 425\,{\scriptsize$\pm$88}\,(5/10) & 136\,{\scriptsize$\pm$2}\,(0/10) & 79\,{\scriptsize$\pm$1}\,(0/10) \\
CEM, mean of five & 460\,{\scriptsize$\pm$17}$^\ast$ & 139\,{\scriptsize$\pm$53}$^\ast$ & 370\,{\scriptsize$\pm$90}\,(3/10) & -- & 73\,{\scriptsize$\pm$1}\,(0/10) \\
CEM, optimistic (UBP2-inspired) & 423\,{\scriptsize$\pm$19} & 195\,{\scriptsize$\pm$43} & 260\,{\scriptsize$\pm$88}\,(2/10) & 129\,{\scriptsize$\pm$2}\,(0/10) & 68\,{\scriptsize$\pm$1}\,(0/10) \\
\bottomrule
\end{tabular}
\end{table*}


\Cref{tab:fixes} and \cref{fig:mech}c compare the fixes at hard search; none wins everywhere. On the walkers, $\chi^2$
weights on the mean of five scorers do best (cheetah 485, hopper 448, against 262 and 356 for standard MPPI), and the
cautious ensemble (mean minus one standard deviation; \citealp{coste2024ensembles}) helps less (295 and 374). On the
drawer with synthetic labels the ranking reverses: the cautious ensemble restores 801 (10/10) while $\chi^2$ weights do
not help (480), consistent with the two patterns above: where the planner reaches states the scorers overrate they
disagree and the penalty finds them; where it chases noise they agree, and only gentler search helps. Automatic temperatures (0.5 times
the disagreement) are the most even (cheetah 521, hopper 569, drawer 650, plate 662); a disagreement guard and the
worst-case-optimal rule with $\kappa=2$ times the disagreement help on the drawer (729, 821) but not on the walkers. With human labels every fix that
searches hard fails (at best 382, 2/10).

Calibrated from 100 pairs of proposals ranked by true return, an oracle check a deployed robot lacks, the rule is the
safest fix (\cref{tab:fixes,tab:calibcl}).
Each of five independent calibration sets per task, drawn on start states the evaluation never uses, calibrates its
own controller: it scores 818--852 on the drawer and 806--1,013 on plate-slide, succeeding in all 50 runs of each,
against 898 and 1,111 for the cloned policy run directly and 611 (8/10) and 351 (6/10) for hard search on the same five
scorers. With human labels every set declines to search (852, uniform averaging). It gains on hopper with every set
(523--550, against 502 for the cloned policy), averages 357 on cheetah (264--445; uniform averaging 233, cloned policy
412) and pays for caution on window-open (309--318, against 392 for standard MPPI and 372).

\emph{Calibrating from people.} Calibrated from one author's comparisons of 100 such pairs, the rule trusts the scorer trained on that author's labels ($\hat\rho=0.63$; $-0.29$ by true return)
and opens the drawer in 5/10 runs (574); the other author's comparisons ($\hat\rho=0.16$) leave it barely searching, at
845 (10/10). Both people called about 70\% of pairs ties, which inflates the concordance on the rest; corrected for
ties, the first author's $\hat\rho$ falls to 0.36 and the rule scores 822 (10/10; \cref{tab:humancalib}). Calibration
against the person who labelled the scorer cannot detect a bias they share; an independent evaluator, the tie
correction or the lower bound limits the damage.

A UBP2-inspired optimistic CEM planner (the CEM component of \citet{nabail2026ubp2} on the five-scorer mean plus one
standard deviation) is the strongest optimizer with the perfect scorer (cheetah 1,244,
plate-slide 1,165) and the most exploited on the arm (plate-slide 68--79 with no successes, drawer 260--425, human
labels 129--136). The collapse persists without optimism (CEM on the plain mean: drawer 370,
plate-slide 73) and at matched budget (MPPI with 768 rollouts: 460, 303), though we do not isolate elite selection.

\paragraph{Labelling where the planner goes.} The scorer is exploited where its labels did not reach, so, as active preference learning
for MPC does \citep{elhasnaouy2026active}, we spend half of a 500-label drawer budget on the planner itself: 250 new
pairs from MPPI episodes planned with a scorer trained on the first 250 offline pairs (\cref{app:setup}). Hard search
then scores 834 (8/10) instead of 504 (11/20), and the cautious ensemble 975 (10/10), level with 2,000 offline labels
and above the cloned policy (898; \cref{fig:mech}a); the distance curve's initial slope turns from $-250$ to $+33$.

\subsection{Best-of-$N$ with a learned verifier}
\label{sec:exp-bon}

As supporting evidence beyond MPPI, we test best-of-$N$ action selection \citep{kwok2025robomonkey} with a state-based
verifier trained on 500 or 2,000 synthetic pairs over perturbed actions (\cref{app:bon}). The $N$ candidates are noisy copies of the cloned action, so
the reference is one noisy action: 755 on the drawer (44/50 successes) and 878 on plate-slide (44/50) on the verifiers'
ten starts, against 898 and 1,111 for the clean action, which a perfect verifier keeps for every $N$ (\cref{fig:kl}c).
With 500 pairs, $N=4$ is level with one noisy action on the drawer within noise (728, though 34/50 successes) and below
it on plate-slide (723, 26/50); from $N=16$ both are well below (530, 375), falling to 388 (14/50) and 143 (0/50) at
$N=1{,}024$. With 2,000 pairs selection beats one noisy action
at every $N$ (drawer 817--895, plate 977--1,061).

\section{Related work}
\label{sec:related}

\paragraph{Reward overoptimization.} \citet{gao2023scaling} fitted how true quality falls with $\sqrt{\KL}$ from the
reference policy under best-of-$n$ and RL against a learned reward model.
Ensembles \citep{coste2024ensembles,eisenstein2024helping,zhai2026uncertainty}, weight averaging \citep{rame2024warm} and
constraints \citep{moskovitz2024confronting} mitigate it. \citet{huang2025chipo} showed that $\chi^2$ regularization gives
guarantees KL regularization cannot, and \citet{huang2025bestofn} analyzed best-of-$N$ and proposed
inference-time pessimism; our \cref{lem:com} centers and tightens their change-of-measure bound, and \cref{thm:chi2}
turns it into a temperature MPPI can check from its own weights and a no-search threshold. \Cref{cor:labels} uses the Bradley--Terry estimation bound of
\citet{zhu2023principled}; \citet{aouad2026revisiting} derive label-count scaling laws for language models. \citet{skalse2022defining,pan2022effects,laidlaw2025correlated} define reward hacking. These works study language
models or RL policies, not a planner in closed loop.

\paragraph{Preferences for robot control.} Preference-based reward learning
\citep{christiano2017deep,sadigh2017active,lee2021pebble,hejna2023fewshot} is known to pick up spurious features
\citep{tien2023causal}; Uni-RLHF \citep{yuan2024unirlhf} provides the crowd labels we use. Learned preference rewards are
now optimized directly by MPC \citep{krupa2025learning,elhasnaouy2026active} and by CEM-based planning in UBP2
\citep{nabail2026ubp2}, which plans optimistically to collect labels; we study such planners at deployment.

\paragraph{Planning against learned models.} Model-based RL guards against exploiting learned dynamics with ensembles,
penalties \citep{chua2018pets,yu2020mopo,kidambi2020morel} or policy constraints
\citep{hansen2024tdmpc2,lin2025tdmpc2policy}; imperfect world models are exploitable \citep{bhamidipaty2026imperfect};
see also conservative MPPI \citep{yoon2026residual}. MPPI is KL-regularized by construction
\citep{williams2017information}; Tsallis-divergence MPPI \citep{wang2021tsallis} gives the $\chi^2$ rule. CriticHack
\citep{luo2026critichack} evaluates visual reward models under policy optimization.

\section{Limitations and conclusion}
\label{sec:conclusion}

Our tasks are simulated, our scorers are small MLPs on state, the human labels come from the two authors on
two tasks, and most cells have 10 runs. Planning uses the true simulator, which isolates the reward but omits
dynamics-model errors. Most calibration pairs are ranked by true return or a
simulated labeller, an oracle a deployed robot lacks, apart from two people's 100 comparisons; the best-of-$N$ verifiers
use synthetic labels. The guarantee of \cref{prop:loop} covers only the selection terms, and
the label-count corollary assumes a linear reward. Window-open, button-press and two Adroit hand tasks are null or limit cases
(\cref{app:null}).

For practitioners, a planner's temperature, sample count and elite fraction are an optimization budget against the
scorer, set by its quality \emph{on the planner's own proposals}. A calibrated temperature was the safest default
tested, given an evaluator independent of the labels; labelling where the planner goes was the most effective fix.

\section*{Impact Statement}
This work shows that robots which optimize rewards learned from human preferences can behave worse than the policy
they start from, without warning from standard validation, and identifies settings and checks that reduce that risk.
The intended effect is safer deployment of learned objectives in robotics. The experiments are in simulation and
involve no human subjects beyond the two authors' own labels.

\section*{Use of Generative AI}
During the preparation of this work the authors used Claude (Anthropic) to search the literature, to write and run
the experiment code, and to draft and edit the text. The authors reviewed and edited the content as needed and take
full responsibility for it. The code includes a script that recomputes every number in the paper from the stored
result files.

\bibliography{refs}
\bibliographystyle{icml2026}

\newpage
\appendix
\onecolumn

\section{Proofs}
\label{app:proofs}

\paragraph{Proof of \cref{lem:com}.} With $w=dp/dq$, $\hat\Delta(p)-\Delta(p)=\E_p[e]-\E_q[e]=\E_q[(w-1)(e-\E_qe)]$, which by
Cauchy--Schwarz under $q$ is at most $\Vert w-1\Vert_{L^2(q)}\,\mathrm{sd}_q(e)=\sigma_e\,s(p)$ in absolute value. For
$p\ne q$, the error $e=\sigma_e(w-1)/s(p)$ has spread $\sigma_e$ and attains the bound. \qed

\paragraph{Proof of \cref{thm:chi2}.} (i) By Cauchy--Schwarz, $\hat\Delta(p)=\E_q[(w-1)(\Jh-\E_q\Jh)]\le s(p)\,\sigma_q(\Jh)$, so
$W_\sigma(p)\le(\sigma_q(\Jh)-\sigma)\,s(p)\le0=W_\sigma(q)$.
(ii) $W_\sigma(p)=\E_q[w\Jh]-\E_q[\Jh]-\sigma\Vert w-1\Vert_{L^2(q)}$ is concave in $w$ on the convex set
$\{w\ge0,\ \E_qw=1\}$, so a point that satisfies the Karush--Kuhn--Tucker conditions is a global maximizer. At $w\ne1$
the gradient of $\Vert w-1\Vert$ is $(w-1)/s$; with a multiplier $\nu$ for $\E_qw=1$ and multipliers $\kappa\ge0$ with
$\kappa w=0$ for $w\ge0$, stationarity reads $\Jh-\sigma(w-1)/s-\nu+\kappa=0$, whose solution is
$w=\max(0,(\Jh-c)/\lambda)$ with $\lambda=\sigma/s$ and $c$ set by normalization: the $\chi^2$ MPPI weights at a
temperature with $\lambda s(p_\lambda)=\sigma$. Conversely, the $\chi^2$ weights at any $\lambda$ satisfy these conditions
with $\sigma=\lambda s(p_\lambda)$, so $p_\lambda$ maximizes $W_{\lambda s(p_\lambda)}$. Next, $p_\lambda$ maximizes
$\E_q[w\Jh]-\tfrac\lambda2\E_q[(w-1)^2]$, so it maximizes $\hat\Delta$ among all $p$ with $s(p)\le s(p_\lambda)$; call
this maximum $g(s)$. $g$ is concave: if $w_1,w_2$ attain $g(s_1),g(s_2)$, the mixture $tw_1+(1-t)w_2$ is feasible at
$ts_1+(1-t)s_2$ by the triangle inequality and has $\hat\Delta=tg(s_1)+(1-t)g(s_2)$. Along the path $W_\sigma=g(s)-\sigma
s$ is therefore concave in $s$, so it rises to its maximizer and falls after, and the maximizer of $g(s)-\sigma s$ moves
to larger $s$ as $\sigma$ falls, so $\lambda s(p_\lambda)$ decreases as $s$ grows. On $N$ samples, for
$\lambda\ge\E_q\Jh-\min\Jh$ no weight is clipped, $w=1+(\Jh-\E_q\Jh)/\lambda$, $s=\sigma_q(\Jh)/\lambda$ and
$\lambda s=\sigma_q(\Jh)$; as $\lambda\to0$, $\lambda s\to0$ because $s\le\sqrt{N-1}$. The product is continuous in
$\lambda$, so $\lambda^\star$ exists for $\sigma>0$; for $\sigma=0$, $W_0=\hat\Delta$ is maximized by putting all weight
on the top-scored plans, the limit of the path as $\lambda\to0$. In both cases and $W_\sigma(p_{\lambda^\star})\ge W_\sigma(p_\lambda)=(\sigma_q(\Jh)-\sigma)s>0$
for any unclipped $\lambda$.
(iii) Unclipped, $w-1=(\Jh-\E_q\Jh)/\lambda$, so $\hat\Delta=\Var_q(\Jh)/\lambda=\sigma_q(\Jh)s$ and
$\Delta=\Cov_q(J,\Jh)/\lambda=\rho_q\sigma_q(J)s$. \qed

\paragraph{Proof of \cref{prop:exec}.} $\bar A_p$ lies in the action box, which is convex. By $L$-smoothness, for every plan $A$,
$J(A)\le J(\bar A_p)+\nabla J(\bar A_p)^\top(A-\bar A_p)+\tfrac L2\Vert A-\bar A_p\Vert^2$. Taking $\E_p$, the linear term
vanishes, which gives the first inequality; the second follows from \cref{lem:com}. For concave $J$, Jensen's inequality
gives $J(\bar A_p)\ge\E_p[J]$ directly. \qed

\paragraph{Proof of \cref{prop:exp}.} As $\lambda\to0$ the exponential weights tend to a point mass on the top plan $i^\star$:
$w=N\mathbb 1_{i^\star}$, $s^2=\E_q[w^2]-1=N-1$, $\hat\Delta=\Jh(i^\star)-\E_q\Jh$. By \cref{lem:com} the worst error is
$e-\E_qe=\sigma(w-1)/\sqrt{N-1}$, which is $\sigma\sqrt{N-1}$ on $i^\star$. If $\max\Jh-\E_q\Jh=o(\sqrt N)$, the guarantee
$\hat\Delta-\sigma\sqrt{N-1}\to-\infty$. By \cref{thm:chi2}, $\max_pW_\sigma\ge W_\sigma(q)=0$. \qed

\begin{corollary}[Labels needed before search is safe]\label{cor:labels}
Let labels follow a correctly specified Bradley--Terry model with $r=\theta^{\star\top}\phi$, $\Vert\phi\Vert\le L$
summed over each segment and $\Vert\theta^\star\Vert\le B$, and let $\Jh$ come from the maximum-likelihood estimate
constrained to $\Vert\theta\Vert\le B$ on $n$ labelled pairs whose feature differences have second-moment matrix
$\Sigma_D$. With $\Sigma_q$ the covariance of plan features under $q$ and the coverage
$C_q=\lambda_{\max}\big((\Sigma_D+\mu I)^{-1/2}\Sigma_q(\Sigma_D+\mu I)^{-1/2}\big)$, with probability $1-\delta$,
$\sigma_e^2\le C^2C_q\big(\tfrac{d+\log(1/\delta)}{\gamma^2n}+\mu B^2\big)$ for any $\mu>0$, with the constants $C$
and $\gamma$ of \citet{zhu2023principled}. A positive gain can be certified once this is below $\sigma_q(\Jh)^2$, which holds for
$n>C^2C_q(d+\log(1/\delta))/\big(\gamma^2(\sigma_q(\Jh)^2-C^2C_q\mu B^2)\big)$ at any $\mu$ with
$C^2C_q\mu B^2<\sigma_q(\Jh)^2$.
\end{corollary}
The labels needed grow linearly with $C_q$, how much further the proposals vary than the labelled pairs did; wider
proposal noise and states the labelled data never reached both raise it.

\paragraph{Proof of \cref{cor:labels}.} The scorer's error on a plan is $e=(\hat\theta-\theta^\star)^\top\Phi$, with $\Phi$ the
plan's summed features, so $\sigma_e^2=v^\top\Sigma_qv$ with $v=\hat\theta-\theta^\star$, which is at most
$C_q\Vert v\Vert^2_{\Sigma_D+\mu I}$. Lemma 3.1 of \citet{zhu2023principled} bounds $\Vert v\Vert_{\Sigma_D+\mu I}$ by
$C\sqrt{(d+\log(1/\delta))/(\gamma^2n)+\mu B^2}$ with probability $1-\delta$, with $\gamma=1/(2+e^{-LB}+e^{LB})$. A positive
gain can be certified when $\sigma_e<\sigma_q(\Jh)$, by \cref{thm:chi2}, and solving the bound for $n$ gives the
condition. \qed

\paragraph{Proof of \cref{prop:onset}.} (i) $\log Z(\eta)=\log\E_q[e^{\eta\Jh}]$ has $\tfrac{d}{d\eta}\log Z=\E_{p_\eta}[\Jh]$.
Then $\tfrac{d}{d\eta}\E_{p_\eta}[J]=\tfrac{d}{d\eta}\E_q[Je^{\eta\Jh}]/Z=\E_{p_\eta}[J\Jh]-\E_{p_\eta}[J]
\E_{p_\eta}[\Jh]$. With $\KL(p_\eta\Vert q)=\eta\E_{p_\eta}[\Jh]-\log Z(\eta)$,
$\tfrac{d}{d\eta}\KL=\E_{p_\eta}[\Jh]+\eta\Var_{p_\eta}(\Jh)-\E_{p_\eta}[\Jh]=\eta\Var_{p_\eta}(\Jh)$. Since $\KL$ is
increasing in $\eta$, the sign of $\tfrac{d}{d\eta}\E_{p_\eta}[J]$ is the sign of the derivative along the distance as
well, and a local maximum needs that sign to change from positive to negative; a covariance that touches zero without
changing sign gives none. (ii) Expanding at $\eta=0$: $\E_{p_\eta}[J]=\E_q[J]+\eta\Cov_q(J,\Jh)+O(\eta^2)$ and
$\KL(p_\eta\Vert q)=\tfrac{\eta^2}{2}\Var_q(\Jh)+O(\eta^3)$, so $\eta=\sqrt2\,d/\sigma_q(\Jh)+O(d^2)$ and
$\E_{p_\eta}[J]-\E_q[J]=\sqrt2\,\Cov_q(J,\Jh)\,d/\sigma_q(\Jh)+O(d^2)=\sqrt2\,\rho_q\sigma_q(J)\,d+O(d^2)$. (iii) The
decomposition is an identity. On $N$ samples with a unique maximizer of $\Jh$, $p_\lambda$ tends to a point mass at
that sample, so $\bar A$ tends to it and $J(\bar A)-\E_p[J]\to0$. \qed

When the probe pools many states, KL and the gain are averaged over states, and the slope in (ii) becomes
$\sqrt2\,\overline{\Cov_q(J,\Jh)}/\sqrt{\overline{\Var_q(\Jh)}}$, which is the predicted slope reported in
Appendix~\ref{app:probe}.

\paragraph{Proof of \cref{prop:loop}.} With deterministic dynamics $Q_t(s_t,a_t)=r_t+V_{t+1}(s_{t+1})$, with $V_T\equiv0$
(and $V\equiv0$ after termination). Summing $Q_t(s_t,a_t)-V_t(s_t)$ over the episode telescopes to
$\sum_tr_t-V_0(s_0)=R-V_0(s_0)$. Adding and subtracting $J(\bar A_t)-J(U_t)$ in each term and splitting it as in
\cref{prop:onset}(iii) gives the identity, along every realization of the proposal noise. If the weights $p_t$ maximize
$W_\sigma$ for some $\sigma\ge\sigma_e(s_t)$, then by \cref{lem:com}
$\Delta_t\ge\hat\Delta_t-\sigma_e(s_t)\,s(p_t)\ge W_\sigma(p_t)\ge W_\sigma(q)=0$. \qed

\section{Setup details}
\label{app:setup}

\paragraph{Planner.} MPPI with $N=256$ sequences, horizon $H=20$, noise $\sigma=0.5$, warm start from the cloned
policy's closed-loop plan, rolled out with the true MuJoCo dynamics (on the arm, a copy of the Meta-World model with
added position sensors whose rollouts match the environment exactly). Temperatures
$\lambda\in\{3,1,0.3,0.1,0.01\}$ on standardized per-step scores. Uniform averaging is MPPI with $\lambda=10^6$:
the weights are uniform and it executes the first action of the average of the 256 clipped proposals. Clipping at the
action bounds makes this average differ from the clean cloned plan $U$ (over 20 steps, by $0.4$ standard deviations of
plan return on the drawer, $0.7$ on cheetah, $1.9$ on plate-slide and $0.03$ on hopper, in the probe), so it is the
$\lambda\to\infty$ end of the same planner, with the same proposals as every other temperature, and the reference for
the distance curves. The cloned policy run directly executes its own action, with no proposals, and scores higher on
every task (cheetah 412, hopper 502, drawer 898, plate-slide 1,111, window-open 372 on the same ten starts).

\paragraph{Statistics.} The resampling unit is a run, one closed-loop episode. A single-scorer cell pairs scorer seeds
0--4 with two or four start states each; an ensemble cell uses one ensemble on ten starts. Standard errors are over
runs, so they include both scorer-seed and start-state variation. Bootstrap intervals resample runs within each cell;
the hierarchical version, used for the slopes and the label-budget trend in \cref{sec:exp-law}, resamples the five
scorer seeds once per draw for all cells of a setting and then runs within each seed. The five calibration sets of \cref{tab:calibcl} are disjoint, so their spread includes the sampling of the
pairs. The automatic temperature sets
$\lambda_t=c\sqrt{\overline{\Var_m\Jh_m(A_i)}}$ from the ensemble's disagreement at each step, $c=0.5$. The guard
compares the ensemble standard deviation of the top 5\% of plans by mean score with its average over all plans and uses
the cautious ensemble with exponential weights when the ratio exceeds 1, otherwise $\chi^2$ weights on the mean. CEM:
6 rounds of 128 samples, top 16 kept, minimum standard deviation 0.05, executing the first action of the final mean. Two variants left out of \cref{tab:fixes} do no better at $\lambda=0.01$: $\chi^2$ weights on the cautious ensemble
(cheetah 439, hopper 348, drawer 625 with 8/10, human labels 229 with 0/10, plate-slide 414 with 7/10) and MPPI on the
optimistic ensemble, mean plus one standard deviation (331 and 221 over 5 runs, 467 with 5/10, 198 with 0/10, 222 with
4/10).

\begin{table}[t]
\centering
\caption{Planner settings. True return of uniform averaging (no search) and of standard MPPI with the learned scorer at $\lambda=3$ and $0.01$, for proposal noise $\sigma$ and horizon $H$; mean $\pm$ standard error over 10 runs (default row: all runs of the main experiments); successes in parentheses.}
\label{tab:sens}
\small
\begin{tabular}{lcccccc}
\toprule
& \multicolumn{3}{c}{Cheetah} & \multicolumn{3}{c}{Drawer, 500 labels} \\
\cmidrule(lr){2-4}\cmidrule(lr){5-7}
& none & $\lambda=3$ & $0.01$ & none & $\lambda=3$ & $0.01$ \\
\midrule
$\sigma=0.5$, $H=20$ (default) & 233\,{\scriptsize$\pm$78} & 480\,{\scriptsize$\pm$38} & 262\,{\scriptsize$\pm$25} & 852\,{\scriptsize$\pm$1} (10/10) & 707\,{\scriptsize$\pm$41} (18/20) & 504\,{\scriptsize$\pm$74} (11/20) \\
$\sigma=0.25$ & 411\,{\scriptsize$\pm$113} & 451\,{\scriptsize$\pm$100} & 517\,{\scriptsize$\pm$39} & 890\,{\scriptsize$\pm$1} (10/10) & 880\,{\scriptsize$\pm$4} (10/10) & 719\,{\scriptsize$\pm$71} (8/10) \\
$\sigma=1$ & 3\,{\scriptsize$\pm$2} & 20\,{\scriptsize$\pm$17} & 3\,{\scriptsize$\pm$24} & 667\,{\scriptsize$\pm$3} (10/10) & 125\,{\scriptsize$\pm$4} (0/10) & 122\,{\scriptsize$\pm$2} (0/10) \\
$H=10$ & 278\,{\scriptsize$\pm$91} & 481\,{\scriptsize$\pm$54} & 328\,{\scriptsize$\pm$36} & 852\,{\scriptsize$\pm$1} (10/10) & 833\,{\scriptsize$\pm$8} (10/10) & 375\,{\scriptsize$\pm$80} (4/10) \\
$H=40$ & 157\,{\scriptsize$\pm$76} & 443\,{\scriptsize$\pm$75} & 173\,{\scriptsize$\pm$23} & 851\,{\scriptsize$\pm$2} (10/10) & 508\,{\scriptsize$\pm$76} (5/10) & 369\,{\scriptsize$\pm$76} (3/10) \\
\bottomrule
\end{tabular}
\end{table}


\paragraph{Scorers and cloned policies.} Scorer: MLP with two hidden layers of 256, Adam (learning rate $3\times10^{-4}$,
weight decay $10^{-4}$), 60 epochs, batch 64, 10\% of pairs held out; output standardized per step on the labelled data.
Cloned policy: MLP with two hidden layers of 256 and tanh output, Adam ($10^{-3}$), 20,000 steps. Locomotion: crowd labels
from Uni-RLHF for the D4RL medium-replay datasets (2,000 pairs, segment length 200), loaded with a restricted unpickler.
Arm: 10 scripted demonstrations (task variations 40--49) for cloning; the labelled data add 100 rollouts of the cloned
policy with action noise 0--1 (variations 10--39); evaluation on variations 0--9.

\paragraph{Labelling where the planner goes.} A scorer trained on the first 250 offline drawer pairs drives standard
MPPI ($\lambda=0.01$ and $0.1$) for 20 episodes on training positions; 250 new pairs compare 50-step segments of those
episodes with each other or with the offline data, labelled by the same synthetic teacher.

\paragraph{Best-of-$N$.} 20,000 state--action-pair comparisons per task from states of the labelled data; both actions
are the cloned action plus $\mathcal N(0,0.5^2)$ noise; the Bradley--Terry teacher prefers the action closer to the
scripted expert action with temperature set by bisection for 85\% agreement. Verifier: MLP on (state, action), two
hidden layers of 256, same optimizer and epochs as the scorers, five seeds per budget. Episodes of 150 steps on
variations 0--9.

\section{One-step probe}
\label{app:probe}

\begin{table*}[t]
\centering
\caption{What the scorer knows about the planner's own proposals. Held-out: label accuracy on held-out pairs and per-step correlation with the true reward on the labelled data. On the proposals, along the cloned policy's episodes (150--200 states): correlation $\rho_q$ between learned and true plan returns over the 256 proposals; the price of the proposal noise $J(U)-\mathbb{E}_qJ$; the selection gain $\Delta$ of standard MPPI at $\lambda=0.01$, on average and at the worst state; and the average gain of the worst-case-optimal rule of \cref{thm:chi2} on the five-scorer mean given the true error spread, whose gain at the worst state is $0.00$ in every row. Gains in units of the proposals' spread $\sigma_q(J)$. Mean $\pm$ standard error over 5 scorer seeds. $\kappa^\star$: median ratio of the true error spread to the ensemble's disagreement.}
\label{tab:probe}
\small
\begin{tabular}{lcccccccc}
\toprule
& \multicolumn{2}{c}{Held-out} & \multicolumn{6}{c}{On the planner's proposals} \\
\cmidrule(lr){2-3}\cmidrule(lr){4-9}
& & & & & \multicolumn{2}{c}{Standard MPPI, $\lambda=0.01$} & \multicolumn{2}{c}{\Cref{thm:chi2}} \\
\cmidrule(lr){6-7}\cmidrule(lr){8-9}
& acc. & corr. & $\rho_q$ & noise & mean $\Delta$ & worst $\Delta$ & mean $\Delta$ & $\kappa^\star$ \\
\midrule
Cheetah & 0.75 & 0.17 & $\phantom{-}0.08$\,{\scriptsize$\pm$0.01} & 1.59 & $\phantom{-}0.18$\,{\scriptsize$\pm$0.07} & $-3.03$\,{\scriptsize$\pm$0.41} & $\phantom{-}0.06$ & 1.4 \\
Hopper & 0.77 & 0.08 & $\phantom{-}0.66$\,{\scriptsize$\pm$0.04} & 0.67 & $\phantom{-}0.22$\,{\scriptsize$\pm$0.19} & $-2.47$\,{\scriptsize$\pm$0.33} & $\phantom{-}0.53$ & 1.2 \\
Drawer 200 & 0.84 & 0.79 & $-0.00$\,{\scriptsize$\pm$0.03} & 0.91 & $-0.64$\,{\scriptsize$\pm$0.18} & $-7.03$\,{\scriptsize$\pm$0.36} & $\phantom{-}0.05$ & 1.8 \\
Drawer 300 & 0.84 & 0.85 & $-0.01$\,{\scriptsize$\pm$0.02} & 0.91 & $-0.21$\,{\scriptsize$\pm$0.06} & $-5.92$\,{\scriptsize$\pm$0.57} & $\phantom{-}0.07$ & 1.9 \\
Drawer 500 & 0.88 & 0.88 & $\phantom{-}0.07$\,{\scriptsize$\pm$0.06} & 0.91 & $\phantom{-}0.16$\,{\scriptsize$\pm$0.12} & $-5.06$\,{\scriptsize$\pm$0.51} & $\phantom{-}0.12$ & 1.4 \\
Drawer 1k & 0.85 & 0.83 & $\phantom{-}0.06$\,{\scriptsize$\pm$0.02} & 0.91 & $\phantom{-}0.12$\,{\scriptsize$\pm$0.08} & $-5.06$\,{\scriptsize$\pm$0.50} & $\phantom{-}0.11$ & 1.9 \\
Drawer 2k & 0.85 & 0.81 & $\phantom{-}0.17$\,{\scriptsize$\pm$0.01} & 0.91 & $\phantom{-}0.43$\,{\scriptsize$\pm$0.07} & $-4.12$\,{\scriptsize$\pm$0.40} & $\phantom{-}0.23$ & 2.0 \\
Drawer human & 0.97 & 0.91 & $-0.17$\,{\scriptsize$\pm$0.02} & 0.91 & $-0.16$\,{\scriptsize$\pm$0.11} & $-6.64$\,{\scriptsize$\pm$0.90} & $\phantom{-}0.05$ & 2.8 \\
Drawer 500, half on planner & 0.84 & 0.92 & $\phantom{-}0.13$\,{\scriptsize$\pm$0.04} & 0.98 & $\phantom{-}0.29$\,{\scriptsize$\pm$0.08} & $-5.66$\,{\scriptsize$\pm$0.45} & $\phantom{-}0.26$ & 1.1 \\
Plate 500 & 0.84 & 0.93 & $\phantom{-}0.42$\,{\scriptsize$\pm$0.03} & 2.16 & $\phantom{-}0.85$\,{\scriptsize$\pm$0.03} & $-4.87$\,{\scriptsize$\pm$1.21} & $\phantom{-}0.42$ & 2.0 \\
Plate 2k & 0.86 & 0.93 & $\phantom{-}0.36$\,{\scriptsize$\pm$0.02} & 2.16 & $\phantom{-}0.98$\,{\scriptsize$\pm$0.10} & $-3.65$\,{\scriptsize$\pm$1.18} & $\phantom{-}0.39$ & 2.0 \\
Window 500 & 0.86 & 0.85 & $\phantom{-}0.36$\,{\scriptsize$\pm$0.02} & 0.33 & $\phantom{-}0.98$\,{\scriptsize$\pm$0.21} & $-8.46$\,{\scriptsize$\pm$2.61} & $\phantom{-}0.39$ & 1.0 \\
Window 2k & 0.85 & 0.91 & $\phantom{-}0.23$\,{\scriptsize$\pm$0.02} & 0.33 & $\phantom{-}0.72$\,{\scriptsize$\pm$0.12} & $-5.68$\,{\scriptsize$\pm$1.07} & $\phantom{-}0.22$ & 2.0 \\
\bottomrule
\end{tabular}
\end{table*}


Along 10 episodes of the cloned policy, at every tenth step, we build the planner's 256 proposals exactly as MPPI does
and roll each out with the true dynamics, recording its learned and true 20-step returns; we also roll out the clean
cloned plan and the weighted-mean plan at 17 temperatures from 100 to 0.01 for both weighting rules. \Cref{tab:probe}
reports the mean correlation $\rho_q$ over states, the price of noise and the selection gain at $\lambda=0.01$ in units
of the mean proposal spread $\sigma_q(J)$, and the code reports the predicted and measured initial slopes and the
distance at which the tilted covariance first reaches zero. For \cref{thm:chi2} we score the proposals with all five
scorers and use their mean $\Jh$. The true error spread is that of $\Jh-J$ after the best positive rescaling of $J$
into the scorer's units, $\sigma_q(\Jh)\sqrt{1-\rho^2}$ (and $\sigma_q(\Jh)$, so no search, when $\rho\le0$); the
rule's temperature solves $\lambda\sqrt{\chi^2}=\sigma$ by bisection. The deployable version replaces $\sigma$ with
$\kappa$ times the ensemble's disagreement, the root mean square over proposals of the standard deviation across
scorers after centering each scorer over the proposals, and $\kappa^\star$ is the per-state ratio of the true error
spread to this disagreement. The calibrated version draws $M$ pairs of proposals at random probed states, labels each
pair by its true plan return, and sets $\hat\rho=\sin(\pi(c-\frac12))$ from the fraction $c$ of pairs the five-scorer
mean orders correctly; the rule then uses $\sigma=\sigma_q(\Jh)\sqrt{1-\hat\rho_+^2}$ at every state, so no search when
$\hat\rho\le0$. The probed states lie on cloned-policy episodes from the ten start states that the closed-loop
evaluation uses. The five calibration sets of \cref{tab:fixes} come instead from a second probe on start states 10--19,
which no closed-loop run uses (on the arm, object positions that the labelled data also cover). The pooled closed-loop runs of \cref{tab:calibcl} use the median $\hat\rho$ over 50 draws of 100 pairs from the first probe (cheetah 0.06, hopper 0.70, drawer 0.06,
human labels $-0.27$, plate-slide 0.49, window-open 0.48), and \cref{app:calib} studies how the estimate varies. The averaging term of \cref{prop:exec},
$J(\bar A)-\E_p[J]$, is $0.25$--$0.85$ standard deviations at uniform weights and shrinks to about zero at hard
exponential search on every task except window-open, where it is $-0.10$ at uniform weights and $-0.15$ at worst, so
there the executed plan can fall short of the guarantee by that much.

\section{Episodes, data support and validation}
\label{app:loop}

\begin{table*}[t]
\centering
\caption{Closed-loop decomposition (\cref{prop:loop}) of standard MPPI episodes against the cloned policy: episode return $R$, the cloned policy's return $V_0$ from the same start, and the episode sums of the selection gain $\Delta_t$, the price of noise $c_t$, the averaging term $\alpha_t$ and the horizon gap $\eta_t$, so that $R-V_0=\sum\Delta-\sum c+\sum\alpha+\sum\eta$ (exact in every run, to $10^{-12}$). $\rho$: mean per-state correlation of the scorer with the true plan return over the 256 proposals, on the states the planner visits (share of states below zero in parentheses) and on the cloned policy's own states (one-step probe, same scorer). Mean over runs, with the standard error over runs for $R-V_0$ and $\sum\Delta$. Each run is paired with $V_0$ from its own start, computed in the planner's simulator copy. On the arm run $k$ uses scorer seed $k$ and start $k$; the walker runs use starts 0 and 1 only (three from start 0, where the cloned cheetah scores 1), so their $V_0$ (cheetah 140, hopper 426) is not the mean over ten starts of \cref{tab:fixes} (412, 502). The cloned cheetah's return also depends sharply on small state differences (from start 1: 347 in the copy, 662 in the environment), the source of its large horizon gap.}
\label{tab:loop}
\small
\setlength{\tabcolsep}{4pt}
\begin{tabular}{lcccccccccc}
\toprule
& runs & $R$ & $V_0$ & $R-V_0$ & $\sum\Delta$ & $\sum c$ & $\sum\alpha$ & $\sum\eta$ & $\rho$ visited & $\rho$ cloned \\
\midrule
Cheetah, $\lambda=0.01$ & 5 & 237 & 140 & $+98$\,{\scriptsize$\pm$68} & $+183$\,{\scriptsize$\pm$153} & $1818$ & $+5$ & $+1727$ & $+0.08$ (34\%) & $+0.05$ \\
Hopper, $\lambda=0.01$ & 5 & 269 & 426 & $-156$\,{\scriptsize$\pm$86} & $-58$\,{\scriptsize$\pm$278} & $754$ & $+1$ & $+655$ & $+0.51$ (13\%) & $+0.76$ \\
Drawer 500, uniform averaging & 5 & 854 & 898 & $-44$\,{\scriptsize$\pm$2} & $0$ & $387$ & $+225$ & $+118$ & $+0.09$ (36\%) & $+0.10$ \\
Drawer 500, $\lambda=3$ & 5 & 765 & 898 & $-133$\,{\scriptsize$\pm$59} & $-31$\,{\scriptsize$\pm$35} & $301$ & $+119$ & $+81$ & $+0.03$ (46\%) & $+0.10$ \\
Drawer 500, $\lambda=0.01$ & 5 & 472 & 898 & $-426$\,{\scriptsize$\pm$164} & $-151$\,{\scriptsize$\pm$77} & $75$ & $0$ & $-200$ & $-0.11$ (62\%) & $+0.10$ \\
Drawer, perfect scorer, $\lambda=0.01$ & 5 & 1099 & 898 & $+200$\,{\scriptsize$\pm$5} & $+632$\,{\scriptsize$\pm$20} & $195$ & $0$ & $-237$ & $+1.00$ (0\%) & -- \\
Drawer human, $\lambda=0.01$ & 5 & 218 & 898 & $-680$\,{\scriptsize$\pm$11} & $-373$\,{\scriptsize$\pm$63} & $69$ & $0$ & $-238$ & $-0.31$ (91\%) & $-0.10$ \\
Plate 500, $\lambda=0.01$ & 5 & 516 & 1115 & $-599$\,{\scriptsize$\pm$108} & $+113$\,{\scriptsize$\pm$58} & $216$ & $+1$ & $-497$ & $-0.07$ (52\%) & $+0.45$ \\
Window 500, $\lambda=0.01$ & 5 & 403 & 381 & $+22$\,{\scriptsize$\pm$17} & $+145$\,{\scriptsize$\pm$42} & $-43$ & $-1$ & $-165$ & $+0.22$ (19\%) & $+0.35$ \\
\bottomrule
\end{tabular}
\end{table*}


\paragraph{Closed-loop decomposition.} \Cref{tab:loop} measures the terms of \cref{prop:loop} along MPPI episodes. At
every visited state the planner's 256 proposals, its executed average and the clean cloned plan are rolled out with the
true reward, and $V_t$ is the true return of the cloned policy from that state to the end of the episode; the identity
then holds in every run to rounding. Uniform averaging pays a price of noise of 387 over a drawer episode and recovers
225 of it by averaging. Hard search removes the averaging term, so it must recover the price by selection: the perfect
scorer does ($+632$), the learned scorers do not on the drawer ($-151$, and $-373$ with human labels). The horizon gap is
large on the walkers ($+655$ to $+1{,}727$), where the cloned policy's own value from a visited state is a poor guide to
what 20-step plans achieve, so the walker rows are best read through $\rho$. On the states the planner visits the
learned scorer's per-state correlation is lower than along the cloned policy's own episodes in every exploited setting
except cheetah, where it is near zero on both; on the drawer with human labels it is negative at 91\% of visited
states.

\paragraph{Data support.} For each arm setting, every step of every labelled segment is a point in the scorer's
standardized input (observation and action). A visited state counts as off the data when its distance to the nearest
labelled point exceeds the 95th percentile of the labelled points' own nearest-neighbour distances (to a point of a
different segment). In observations alone, hard search ($\lambda=0.01$, five recorded episodes) is off the data at 53\%
of steps on the drawer with 500 labels, 60\% with human labels and 51\% on plate-slide, and uniform averaging at none;
on window-open the shares are 75\% and 60\%. With actions included the shares are 42\%, 47\%, 27\% and 32\% for hard
search, and 0--2\% without search. The window-open evaluation positions lie outside the labelled data in the
observation space, and search still helps there, so leaving the data is not by itself what makes a scorer exploitable.
In hand position alone, hard search ends on average 6.6\,cm from the nearest labelled hand position on the drawer with
synthetic labels (0.3\,cm without search) but within 0.8\,cm on plate-slide.

\paragraph{Held-out accuracy without shared trajectories.} The scorers hold out a random 10\% of pairs, so a held-out
pair can share trajectories with training pairs. Splitting the labelled episodes 70/30 instead, and holding out the
pairs whose two segments both come from held-out episodes, gives the same picture (five seeds pooled; a random split
with the same number of training pairs in parentheses): drawer 500 labels 180 of 213 correct, 0.85 (0.85); 2,000 labels
0.83 (0.83); human labels 114 of 118, 0.97 (0.99); plate-slide 0.86 (0.81); window-open 0.84 (0.89); hopper 0.70
(0.74); cheetah 0.72 (0.73).

\section{Calibrating from few pairs}
\label{app:calib}

\begin{table*}[t]
\centering
\caption{Calibrating the error spread from $M$ labelled pairs of the planner's own proposals (one-step probe, 200 random draws of the pairs). $\rho_q$: mean per-state correlation of the five-scorer mean with the true plan return; $\hat\rho$ all: the Gaussian estimate $\sin(\pi(c-\frac12))$ from every pair at every state; below: share of states whose own correlation is under it. Range: 10th and 90th percentile of $\hat\rho$ over draws. Gains $\Delta$ (mean over states) and worst-state $\Delta$ (mean over draws) of the calibrated rule with 100 pairs, in units of $\sigma_q(J)$: with each state's true $\rho$ (worst state $0.00$ in every row), with the point estimate, with the one-sided 95\% lower bound (LCB), and with labels from a noisy teacher that agrees with the true order on 85\% or 70\% of proposal pairs (point estimate).}
\label{tab:calib}
\scriptsize
\setlength{\tabcolsep}{3pt}
\resizebox{\textwidth}{!}{%
\begin{tabular}{lccccccccccccc}
\toprule
& & & & \multicolumn{3}{c}{$\hat\rho$ range, true labels} & true $\rho$ & \multicolumn{2}{c}{point, $M=100$} & \multicolumn{2}{c}{LCB, $M=100$} & 85\% teacher & 70\% teacher \\
\cmidrule(lr){5-7}\cmidrule(lr){8-8}\cmidrule(lr){9-10}\cmidrule(lr){11-12}\cmidrule(lr){13-13}\cmidrule(lr){14-14}
& $\rho_q$ & $\hat\rho$ all & below & $M=25$ & $M=100$ & $M=200$ & $\Delta$ & $\Delta$ & worst & $\Delta$ & worst & $\hat\rho$ / $\Delta$ & $\hat\rho$ / $\Delta$ \\
\midrule
Cheetah & $0.08$ & $0.07$ & 54\% & $-0.31$, $0.43$ & $-0.09$, $0.28$ & $-0.03$, $0.22$ & $0.06$ & $0.03$ & $-0.06$ & $0.00$ & $-0.01$ & $0.06$ / $0.03$ & $0.06$ / $0.03$ \\
Hopper & $0.73$ & $0.69$ & 26\% & $0.43$, $0.88$ & $0.54$, $0.81$ & $0.60$, $0.78$ & $0.53$ & $0.49$ & $-0.02$ & $0.40$ & $-0.01$ & $0.64$ / $0.46$ & $0.43$ / $0.37$ \\
Drawer 200 & $0.01$ & $0.02$ & 46\% & $-0.43$, $0.43$ & $-0.19$, $0.25$ & $-0.13$, $0.16$ & $0.05$ & $0.00$ & $-0.21$ & $0.00$ & $-0.03$ & $0.00$ / $0.00$ & $0.00$ / $0.00$ \\
Drawer 300 & $-0.01$ & $-0.02$ & 45\% & $-0.54$, $0.31$ & $-0.22$, $0.16$ & $-0.16$, $0.11$ & $0.07$ & $0.01$ & $-0.11$ & $0.00$ & $-0.01$ & $0.00$ / $0.01$ & $0.03$ / $0.01$ \\
Drawer 500 & $0.08$ & $0.07$ & 47\% & $-0.31$, $0.43$ & $-0.09$, $0.28$ & $-0.06$, $0.19$ & $0.12$ & $0.04$ & $-0.22$ & $0.01$ & $-0.04$ & $0.06$ / $0.03$ & $0.06$ / $0.03$ \\
Drawer 1k & $0.07$ & $0.07$ & 48\% & $-0.31$, $0.43$ & $-0.13$, $0.25$ & $-0.08$, $0.20$ & $0.11$ & $0.05$ & $-0.15$ & $0.01$ & $-0.02$ & $0.09$ / $0.05$ & $0.09$ / $0.05$ \\
Drawer 2k & $0.19$ & $0.20$ & 46\% & $-0.19$, $0.54$ & $0.03$, $0.43$ & $0.06$, $0.32$ & $0.23$ & $0.14$ & $-0.15$ & $0.05$ & $-0.05$ & $0.22$ / $0.14$ & $0.13$ / $0.11$ \\
Drawer human & $-0.18$ & $-0.21$ & 53\% & $-0.64$, $0.19$ & $-0.40$, $-0.03$ & $-0.34$, $-0.09$ & $0.05$ & $0.00$ & $-0.02$ & $0.00$ & $0.00$ & $-0.13$ / $0.00$ & $-0.06$ / $-0.01$ \\
Drawer, half on planner & $0.16$ & $0.18$ & 53\% & $-0.31$, $0.54$ & $0.00$, $0.37$ & $0.02$, $0.31$ & $0.26$ & $0.12$ & $-0.17$ & $0.03$ & $-0.04$ & $0.13$ / $0.11$ & $0.13$ / $0.10$ \\
Plate 500 & $0.44$ & $0.45$ & 19\% & $0.06$, $0.73$ & $0.25$, $0.61$ & $0.32$, $0.58$ & $0.42$ & $0.35$ & $-0.03$ & $0.23$ & $-0.01$ & $0.43$ / $0.34$ & $0.28$ / $0.27$ \\
Plate 2k & $0.39$ & $0.40$ & 21\% & $0.05$, $0.73$ & $0.22$, $0.56$ & $0.28$, $0.54$ & $0.39$ & $0.30$ & $0.00$ & $0.19$ & $0.00$ & $0.37$ / $0.29$ & $0.28$ / $0.24$ \\
Window 500 & $0.44$ & $0.47$ & 47\% & $0.06$, $0.81$ & $0.28$, $0.61$ & $0.38$, $0.59$ & $0.39$ & $0.33$ & $-0.03$ & $0.22$ & $-0.01$ & $0.43$ / $0.31$ & $0.28$ / $0.24$ \\
Window 2k & $0.24$ & $0.28$ & 59\% & $-0.19$, $0.64$ & $0.06$, $0.45$ & $0.14$, $0.38$ & $0.22$ & $0.12$ & $-0.86$ & $0.05$ & $-0.40$ & $0.19$ / $0.10$ & $0.09$ / $0.07$ \\
\bottomrule
\end{tabular}}
\end{table*}


\Cref{tab:calib} repeats the calibration of \cref{app:probe} with 200 random draws of exactly $M$ pairs (pairs tied in
true return are redrawn) for $M=25$, 50, 100 and 200, and three labellers: the true order, and Bradley--Terry teachers
$P(i\succ j)=1/(1+e^{-(J_i-J_j)/\tau})$ with $\tau$ set so that they agree with the true order on 85\% or 70\% of
proposal pairs, a stand-in for a person comparing two short plans from the same state. A cautious version replaces $c$
by its one-sided Wilson lower bound at nominal 95\%.

\emph{Sampling noise.} The 10--90\% range of $\hat\rho$ spans 0.45--0.85 with 25 pairs, 0.27--0.44 with 100 and
0.18--0.29 with 200. With 100 pairs the point estimate loses at most 0.22 standard deviations at the worst state on
average over draws in twelve of the thirteen settings, but up to 0.47 in the unluckiest tenth of draws (drawer, 200
labels). \emph{Lower bound.} The lower bound cuts the worst-state loss to at most 0.05 (0.40 on window-open with 2,000
labels) and the unluckiest-tenth loss to 0.25, at a cost in gain: it keeps 77\% of the true-$\rho$ rule's gain on
hopper, 48--57\% on plate-slide and on window-open with 500 labels, and 12--24\% in the other settings where that gain
exceeds 0.2. \emph{Noisy labels.} A noisy teacher shrinks $\hat\rho$ toward zero (hopper 0.71, 0.64 and 0.43 for the
true order, 85\% and 70\%; plate-slide 0.45, 0.43 and 0.28), so where the scorer is informative the rule searches more
gently: with 100 pairs the worst-state loss stays at most 0.19 (85\%) and 0.18 (70\%) in the same twelve settings, and
the share of the gain kept falls from 47--93\% to 42--87\% and 31--71\%. \emph{Gaussian assumption.} Over all pairs,
$\hat\rho$ is within 0.05 of the mean Pearson correlation in every setting, and above it by up to 0.04 on plate-slide
and window-open, where it is mildly optimistic. The larger approximation is pooling: one $\hat\rho$ serves every state,
and 19--59\% of states have their own correlation below it. Those states cause the losses, and window-open with 2,000
labels has the most of them (59\%).

\begin{table*}[t]
\centering
\caption{The calibrated rule in closed loop on the same ten starts per column. Top: references, and hard search ($\lambda=0.01$) with single scorers (20 runs) and on the five-scorer mean that the calibrated rule uses. Pooled: the median $\hat\rho$ over 50 draws of 100 proposal pairs ranked by true return (point estimate and one-sided Wilson lower bound at nominal 95\%) or by a noisy teacher that agrees with the true order on 85\% of pairs. Independent sets: five disjoint sets of exactly 100 pairs, each calibrating its own controller; mean $\pm$ standard deviation across sets, then the worst set. Held-out sets: the same, with the pairs drawn on start states 10--19, which the evaluation never uses. $\kappa=\sqrt{1-\hat\rho^2}$ is the error spread as a share of the scorer's spread over the proposals, and $\kappa=1$ means no search (uniform averaging). Otherwise mean $\pm$ standard error over runs; successes over runs in parentheses.}
\label{tab:calibcl}
\small
\setlength{\tabcolsep}{4pt}
\resizebox{\textwidth}{!}{%
\begin{tabular}{lcccccc}
\toprule
& Cheetah & Hopper & Drawer 500 & Drawer human & Plate 500 & Window 500 \\
\midrule
Cloned policy, run directly & 412\,{\scriptsize$\pm$114} & 502\,{\scriptsize$\pm$32} & 898\,{\scriptsize$\pm$1}\,(10/10) & 898\,{\scriptsize$\pm$1}\,(10/10) & 1111\,{\scriptsize$\pm$5}\,(10/10) & 372\,{\scriptsize$\pm$13}\,(10/10) \\
Uniform averaging (no search) & 233\,{\scriptsize$\pm$78} & 418\,{\scriptsize$\pm$38} & 852\,{\scriptsize$\pm$1}\,(10/10) & 852\,{\scriptsize$\pm$1}\,(10/10) & 1051\,{\scriptsize$\pm$6}\,(10/10) & 295\,{\scriptsize$\pm$9}\,(10/10) \\
Standard MPPI, single scorers & 262\,{\scriptsize$\pm$25} & 356\,{\scriptsize$\pm$39} & 504\,{\scriptsize$\pm$74}\,(11/20) & 216\,{\scriptsize$\pm$10}\,(0/20) & 436\,{\scriptsize$\pm$65}\,(15/20) & 392\,{\scriptsize$\pm$14}\,(10/10) \\
Standard MPPI, five-scorer mean & 296\,{\scriptsize$\pm$33} & 339\,{\scriptsize$\pm$66} & 611\,{\scriptsize$\pm$89}\,(8/10) & 206\,{\scriptsize$\pm$13}\,(0/10) & 351\,{\scriptsize$\pm$86}\,(6/10) & 425\,{\scriptsize$\pm$16}\,(10/10) \\
\midrule
Pooled, true order, point & 451\,{\scriptsize$\pm$54} & 538\,{\scriptsize$\pm$30} & 835\,{\scriptsize$\pm$6}\,(10/10) & 852\,{\scriptsize$\pm$1}\,(10/10) & 987\,{\scriptsize$\pm$18}\,(10/10) & 311\,{\scriptsize$\pm$9}\,(10/10) \\
\quad $\kappa$ & 0.998 & 0.713 & 0.998 & 1.000 & 0.872 & 0.876 \\
Pooled, true order, lower bound & 264\,{\scriptsize$\pm$88} & 519\,{\scriptsize$\pm$35} & 852\,{\scriptsize$\pm$1}\,(10/10) & 852\,{\scriptsize$\pm$1}\,(10/10) & 1036\,{\scriptsize$\pm$6}\,(10/10) & 306\,{\scriptsize$\pm$9}\,(10/10) \\
\quad $\kappa$ & 1.000 & 0.856 & 1.000 & 1.000 & 0.977 & 0.970 \\
Pooled, 85\% teacher, point & 451\,{\scriptsize$\pm$54} & 532\,{\scriptsize$\pm$31} & 835\,{\scriptsize$\pm$6}\,(10/10) & 852\,{\scriptsize$\pm$1}\,(10/10) & 1007\,{\scriptsize$\pm$13}\,(10/10) & 310\,{\scriptsize$\pm$9}\,(10/10) \\
\quad $\kappa$ & 0.998 & 0.771 & 0.998 & 1.000 & 0.905 & 0.905 \\
\midrule
Five independent sets, point & 307\,{\scriptsize$\pm$52}, 264 & 551\,{\scriptsize$\pm$23}, 517 & 838\,{\scriptsize$\pm$13}, 824\,(50/50) & 851\,{\scriptsize$\pm$3}, 846\,(50/50) & 1009\,{\scriptsize$\pm$6}, 999\,(50/50) & 312\,{\scriptsize$\pm$5}, 308\,(50/50) \\
\quad $\kappa$ & 0.976--1.000 & 0.588--0.729 & 0.969--1.000 & 0.996--1.000 & 0.891--0.918 & 0.771--0.941 \\
Five independent sets, lower bound & 264\,{\scriptsize$\pm$0}, 264 & 538\,{\scriptsize$\pm$22}, 500 & 852\,{\scriptsize$\pm$0}, 852\,(50/50) & 852\,{\scriptsize$\pm$0}, 852\,(50/50) & 1039\,{\scriptsize$\pm$2}, 1036\,(50/50) & 306\,{\scriptsize$\pm$2}, 303\,(50/50) \\
\quad $\kappa$ & 1.000 & 0.757--0.873 & 1.000 & 1.000 & 0.977--0.989 & 0.903--0.996 \\
Five held-out sets, point & 357\,{\scriptsize$\pm$86}, 264 & 539\,{\scriptsize$\pm$12}, 523 & 840\,{\scriptsize$\pm$16}, 818\,(50/50) & 852\,{\scriptsize$\pm$0}, 852\,(50/50) & 959\,{\scriptsize$\pm$86}, 806\,(50/50) & 313\,{\scriptsize$\pm$3}, 309\,(50/50) \\
\quad $\kappa$ & 0.960--1.000 & 0.562--0.771 & 0.951--1.000 & 1.000 & 0.750--0.918 & 0.809--0.918 \\
Five held-out sets, lower bound & 264\,{\scriptsize$\pm$0}, 264 & 509\,{\scriptsize$\pm$30}, 472 & 848\,{\scriptsize$\pm$8}, 835\,(50/50) & 852\,{\scriptsize$\pm$0}, 852\,(50/50) & 1028\,{\scriptsize$\pm$17}, 998\,(50/50) & 306\,{\scriptsize$\pm$2}, 304\,(50/50) \\
\quad $\kappa$ & 1.000 & 0.734--0.903 & 0.998--1.000 & 1.000 & 0.889--0.989 & 0.930--0.989 \\
\bottomrule
\end{tabular}}
\end{table*}


\emph{Closed loop.} \Cref{tab:calibcl} runs the rule with the pooled estimate, its lower bound and a noisy teacher, and
with five independent sets of exactly 100 pairs per setting (500 distinct pairs, none tied in true return), each
calibrating its own controller on the same ten starts, drawn either from the probe on the evaluation's start states or
from a second probe on start states 10--19 that the evaluation never uses. The noisy teacher changes little (within 20
points of the true-order version). Across independent sets the arm results barely move (evaluation starts: drawer
824--852, plate-slide 999--1,013, window-open 308--320; held-out starts: 818--852, 806--1,013, 309--318; every run
successful); the walkers move more (cheetah 264--392 and 264--445, hopper 517--575 and 523--550), since a small change
in $\kappa$ near 1 changes how hard they search. On plate-slide one held-out set estimates $\hat\rho=0.66$, against
0.40--0.51 for the others, so it searches harder and scores 806, still succeeding in every run. With the lower bound every set keeps the
arm at or near uniform averaging's level (evaluation starts: drawer 852 with either label set, plate-slide
1,036--1,041, window-open 303--310; held-out starts: drawer 835--852, plate-slide 998--1,041, window-open 304--308;
every run successful), declines to search on cheetah (264 with every set), and keeps most of hopper's gain (500--560
and 472--551).

What the estimate must get right is stated by the following consequence of \cref{lem:com} and \cref{thm:chi2}.

\begin{corollary}[Calibrated rule]\label{cor:calib}
(i) At a state whose proposals have $\rho_s=\mathrm{Corr}_q(J,\Jh)>0$, the rule with
$\sigma=\sigma_q(\Jh)\sqrt{1-\hat\rho_+^2}$, $\hat\rho_+=\max(0,\hat\rho)$, does not lose ($\Delta\ge0$) whenever
$\hat\rho\le\rho_s$, and it does not search when $\hat\rho\le0$.
\\[2pt](ii) Let a labeller rank pairs by $J+u$, with $J$ in the scorer's units, $\Jh=J+e$, and $e$ and $u$ uncorrelated
with $J$. If $(\Jh,J+u)$ is bivariate Gaussian under $q$, the population concordance $c$ of the labeller with $\Jh$ on
two independent proposals satisfies $\sin(\pi(c-\frac12))=\rho_u=\mathrm{Corr}_q(\Jh,J+u)$, and $\rho_u>\rho_q$
exactly when $\Cov_q(e,u)>\sigma_q(J)\big(\sqrt{\sigma_q(J)^2+\sigma_u^2}-\sigma_q(J)\big)$, about $\sigma_u^2/2$
for a small labelling error; with $\Cov_q(e,u)=0$, $\rho_u<\rho_q$ whenever $\sigma_u>0$.
\end{corollary}

\emph{Proof.} (i) By the scale argument $\sigma_e=\sigma_q(\Jh)\sqrt{1-\rho_s^2}\le\sigma$, so the rule's weights $p$
maximize $W_\sigma$ and $\Delta(p)\ge\hat\Delta(p)-\sigma_es(p)\ge W_\sigma(p)\ge W_\sigma(q)=0$ by \cref{lem:com};
if $\hat\rho\le0$ then $\hat\rho_+=0$, $\sigma=\sigma_q(\Jh)$ and \cref{thm:chi2}(i) gives $p=q$. (ii) For a bivariate
Gaussian, two independent draws are concordant with probability $\frac12+\arcsin(\rho_u)/\pi$
\citep{kruskal1958ordinal}; and $\rho_u=\big(\sigma_q(J)^2+\Cov_q(e,u)\big)\big/\big(\sigma_q(\Jh)\sqrt{\sigma_q(J)^2+\sigma_u^2}\big)$ and
$\rho_q=\sigma_q(J)/\sigma_q(\Jh)$. \qed

In the population, then, a labeller whose mistakes are independent of the scorer's only lowers the correlation that
$\hat\rho$ estimates, as the noisy teachers do; one whose mistakes the scorer has learned, because it was trained on
that labeller's labels, can raise it above the true correlation, and then the rule searches where it should not. A
finite, pooled estimate can also exceed a given state's correlation through sampling noise (the ranges above) and
pooling over states. The one-sided Wilson bound at nominal 95\% makes $\hat\rho$ at most the pooled correlation with
probability close to 0.95 under the Gaussian model (its coverage is approximate: 94.2\% for 100 independent pairs at
90\% agreement), so (i) then covers every state whose own correlation is at least the pooled one. This is a statement about the pooled correlation, not a simultaneous one over states: the states
it leaves uncovered are the 19--59\% noted above, so the rule is an empirical safeguard, not a guarantee.

\emph{Ties.} The true-order and teacher labels above never tie. People often cannot tell two plans apart, and $c$ is
then measured on the decided pairs only. If a person calls a tie when the two plans' perceived returns differ by less
than a threshold, as in threshold models of paired comparison \citep{rao1967ties}, the decided pairs are the clearer
ones and agree with the scorer more often than random pairs do, so $\sin(\pi(c-\frac12))$ overstates $\rho_u$. In the
Gaussian model of \cref{cor:calib}(ii), let the labeller's difference $X$ (standardized) and the scorer's difference have
correlation $\rho_u$, and let a share $t$ of pairs be tied, those with $|X|<z_t=\Phi^{-1}(\frac{1+t}2)$. Given $X$, the
scorer agrees in sign with probability $\Phi(\rho_u|X|/\sqrt{1-\rho_u^2})$, so the decided pairs' concordance is
\[c_t(\rho_u)=\E\big[\Phi\big(\rho_u|X|/\sqrt{1-\rho_u^2}\big)\,\big|\,|X|>z_t\big],\]
which is increasing in $\rho_u$, equals $\frac12+\arcsin(\rho_u)/\pi$ at $t=0$, and grows with $t$ when $\rho_u>0$. We
estimate $\rho_u$ by solving $c_t(\hat\rho)=c$ at the observed share of ties, and its lower bound by applying the same
map to the Wilson bound on $c$. On the probe's proposals, a true-order labeller that ties the closest 70\% of pairs gives
an uncorrected $\hat\rho$ of 0.97 on hopper, 0.77 on plate-slide and 0.73 on window-open, against 0.69, 0.45 and 0.47
over all pairs; the corrected estimates are 0.70, 0.45 and 0.42, and within 0.03 on cheetah and the drawer with 500
labels. With human labels, where the correlation is negative ($-0.21$ over all
pairs), both estimates keep the sign but shrink it ($-0.08$ and $-0.04$), so the rule still does not search.

\emph{Two people's comparisons.} One author compared 100 random pairs of the drawer proposals blind: the 20 planned
steps of two proposals from the same state, as 2-s clips in random order and sides (9 minutes in all). They could not
tell 68 apart and agreed with the true return on 16 of the other 32, which is chance. Where one plan opened the drawer
more than 1\,mm further over the clip (14 decided pairs) they chose it 13 times; the true 20-step return chose it 6
times, and mostly favours the plan that ends with the hand nearer the handle (68\% of 31 such pairs, against 42\% for
the person). The scorers with 500 and 2,000 synthetic labels agree with these choices at chance ($\hat\rho=0.00$, lower
bound $-0.42$), so the rule does not search with them and keeps 852 (10 of 10). The scorer trained on the same
person's clip labels agrees with them ($\hat\rho=0.63$, lower bound 0.23) although it disagrees with the true return on
the same pairs ($-0.29$). Calibrated from these choices the rule searches with it ($\kappa=0.77$) and scores 574 (5 of
10 opened), against 852 when calibrated by true return; with the lower bound ($\kappa=0.97$) it scores 839 (10 of 10).
Two thirds of the pairs were ties, so the decided pairs are the clearer ones; corrected for ties, $\hat\rho=0.36$
(lower bound 0.12) and the rule searches gently ($\kappa=0.93$), scoring 822 (10 of 10), still below uniform
averaging: the correction removes the overconfidence the ties add, not the shared bias.
The other author compared the same 100 pairs (38 minutes): they could not tell 71 apart and agreed with the true return
on 12 of the other 29, and where both people decided (12 pairs) they agreed 9 times. Their choices do not back the
scorer trained on the first author's clip labels ($\hat\rho=0.16$, lower bound $-0.30$, so $\kappa=0.99$): the rule
barely searches and scores 845 (10 of 10). Pooled over both people, $\hat\rho=0.42$ (lower bound 0.11) gives 812 and
846 (10 of 10 each). On the scorers with 500 and 2,000 synthetic labels the second author's choices give
$\hat\rho=0.16$ and 0.37, and the rule scores 831 and 893 (10 of 10 each). \Cref{tab:humancalib} collects the three kinds of calibration
for this scorer, with and without the tie correction. They are not interchangeable: the oracle sees that the scorer
ranks the proposals worse than chance and declines to search; the labeller who trained the scorer shares its bias and
endorses it; the other author does not. A rule calibrated from a person's preferences checks the scorer against that
person, so it cannot catch a scorer that learned a bias it shares with the evaluator (\cref{cor:calib}(ii)); a second
person, the tie correction or the lower bound limits the damage. With two labellers on one task, this is a case study,
not an estimate of how often it happens.

\begin{table}[t]
\centering
\caption{The calibrated rule on the drawer scorer trained on one author's clip labels ($\rho_q=-0.18$ on the planner's proposals, \cref{tab:calib}), by who ranks the 100 calibration pairs. $\hat\rho$ with its one-sided Wilson lower bound in parentheses, as measured on the decided pairs and corrected for can't-tell answers; then the true return (successes over 10 runs) of the rule at each estimate. Uniform averaging scores 852 (10/10) and the cloned policy 898.}
\label{tab:humancalib}
\small
\setlength{\tabcolsep}{3pt}
\resizebox{\textwidth}{!}{%
\begin{tabular}{lccccccc}
\toprule
& decided & \multicolumn{2}{c}{$\hat\rho$ (lower bound)} & \multicolumn{2}{c}{return, as measured} & \multicolumn{2}{c}{return, ties corrected} \\
\cmidrule(lr){3-4}\cmidrule(lr){5-6}\cmidrule(lr){7-8}
Calibration pairs ranked by & pairs & as measured & ties corrected & point & lower bound & point & lower bound \\
\midrule
Oracle (true return), five sets & 100 of 100 & $-0.45$ to $-0.09$ & (no ties) & 852 (10/10) & 852 (10/10) & -- & -- \\
Author whose labels trained the scorer & 32 of 100 & $0.63$ ($0.23$) & $0.36$ ($0.12$) & 574 (5/10) & 839 (10/10) & 822 (10/10) & 846 (10/10) \\
The other author & 29 of 100 & $0.16$ ($-0.30$) & $0.08$ ($-0.16$) & 845 (10/10) & 852 (10/10) & 846 (10/10) & 852 (10/10) \\
Both authors pooled & 61 of 200 & $0.42$ ($0.11$) & $0.23$ ($0.06$) & 812 (10/10) & 846 (10/10) & 838 (10/10) & 846 (10/10) \\
\bottomrule
\end{tabular}}
\end{table}


\section{Null and limit cases}
\label{app:null}

\paragraph{Button-press.} Same protocol as the drawer, 5 runs per setting: not exploited at any budget (standard MPPI
151--157 at $\lambda=3$ and 154--171 at $\lambda=0.01$ with 200, 500 or 2,000 labels; nearly every run presses the
button). Here the perfect scorer does worse (126--140), because over a 20-step horizon Meta-World's dense reward pays
for hovering near the button rather than pressing it, while the preference-trained scorers learned progress over 50-step
clips. The label budget at which exploitation starts is task-specific.

\paragraph{Adroit door and hammer.} With crowd labels that compare only successful human demonstrations
\citep{rajeswaran2018learning,yuan2024unirlhf}, the scorer never opens the door (0 of 12 runs at any temperature) and
drives the nail in 2 of 60 runs, while the perfect scorer opens the door in 6 of 6 runs at $\lambda\le0.01$ and drives
the nail in 4 of 5 at $\lambda=0.1$. With no failures in the labelled data the scorer has no signal about failure, and
no weighting rule can supply one.

\section{Best-of-$N$ details}
\label{app:bon}
At each step the cloned policy's action plus $\mathcal N(0,0.5^2)$ noise gives $N$ candidates, and an MLP verifier
$V(s,a)$ picks one. It is trained with the Bradley--Terry loss on pairs of perturbed actions, where a synthetic labeller
prefers the action closer to Meta-World's scripted expert, with noise set for 85\% agreement. The verifier sees the
state, not images, and its labels are synthetic. The perfect verifier is the negative distance to the expert action.
Single verifiers pool scorer seeds 0--4 with episode seeds 0--9 (50 runs per $N$); the perfect verifier and one
randomly perturbed action, with no verifier, use the same ten starts with five noise draws each. In \cref{fig:kl}c,
$N=1$ is one noisy action and the clean cloned action, executed without noise, is the reference line. A cautious five-verifier ensemble does not help on
plate-slide (0 of 10 at $N\ge16$). The verifiers' held-out pair accuracies are 0.62--0.80 at 500 pairs and 0.69--0.81
at 2,000, again a poor guide to how they behave under selection. Best-of-$N$ is itself an exponential tilt of $q$, the action noise. For continuous scores the
population best-of-$N$ choice has density $dp_N/dq=N\,F(\Jh)^{N-1}$, with $F$ the distribution function of $\Jh$ under
$q$: this is $q$ tilted by $e^{\eta\log F(\Jh)}$ at $\eta=N-1$, so \cref{prop:onset}(i) applies with $\log F(\Jh)$ in
place of $\Jh$, and true return rises with $N$ exactly while $\Cov_{p_N}(J,\log F(\Jh))>0$. Its KL to $q$ is
$\log N-(N-1)/N$ \citep{stiennon2020learning,beirami2024theoretical}, which puts $N=4$ to $1{,}024$ at $d=0.8$ to $2.4$,
the range of \cref{fig:kl}a. This is the population value; the selected index on a finite set of $N$ candidates has KL
$\log N$ to a uniform pick.

\end{document}
